\documentclass[11pt,a4paper]{article}
\usepackage[T1]{fontenc}
\usepackage{lmodern}
\pdfgentounicode=1
\pdfglyphtounicode{tfm:lmsy10/mapsto}{007C}
\pdfglyphtounicode{tfm:lmex10/parenleftbig}{0028}
\pdfglyphtounicode{tfm:lmex10/parenrightbig}{0029}
\pdfglyphtounicode{tfm:lmex10/bracketleftbig}{005B}
\pdfglyphtounicode{tfm:lmex10/bracketrightbig}{005D}
\pdfglyphtounicode{tfm:lmex10/vextendsingle}{23D0}
\pdfglyphtounicode{tfm:lmex10/parenleftBig}{0028}
\pdfglyphtounicode{tfm:lmex10/parenrightBig}{0029}
\pdfglyphtounicode{tfm:lmex10/parenleftbigg}{0028}
\pdfglyphtounicode{tfm:lmex10/parenrightbigg}{0029}
\pdfglyphtounicode{tfm:lmex10/bracketleftbigg}{005B}
\pdfglyphtounicode{tfm:lmex10/bracketrightbigg}{005D}
\pdfglyphtounicode{tfm:lmex10/ceilingleftbigg}{2308}
\pdfglyphtounicode{tfm:lmex10/ceilingrightbigg}{2309}
\pdfglyphtounicode{tfm:lmex10/parenleftBigg}{0028}
\pdfglyphtounicode{tfm:lmex10/parenrightBigg}{0029}
\pdfglyphtounicode{tfm:lmex10/bracketleftBigg}{005B}
\pdfglyphtounicode{tfm:lmex10/bracketrightBigg}{005D}
\pdfglyphtounicode{tfm:lmex10/braceleftBigg}{007B}
\pdfglyphtounicode{tfm:lmex10/bracerightBigg}{007D}
\pdfglyphtounicode{tfm:lmex10/parenlefttp}{239B}
\pdfglyphtounicode{tfm:lmex10/parenrighttp}{239E}
\pdfglyphtounicode{tfm:lmex10/bracketlefttp}{23A1}
\pdfglyphtounicode{tfm:lmex10/bracketrighttp}{23A4}
\pdfglyphtounicode{tfm:lmex10/bracketleftbt}{23A3}
\pdfglyphtounicode{tfm:lmex10/bracketrightbt}{23A6}
\pdfglyphtounicode{tfm:lmex10/parenleftbt}{239D}
\pdfglyphtounicode{tfm:lmex10/parenrightbt}{23A0}
\pdfglyphtounicode{tfm:lmex10/summationtext}{2211}
\pdfglyphtounicode{tfm:lmex10/producttext}{220F}
\pdfglyphtounicode{tfm:lmex10/integraltext}{222B}
\pdfglyphtounicode{tfm:lmex10/summationdisplay}{2211}
\pdfglyphtounicode{tfm:lmex10/productdisplay}{220F}
\pdfglyphtounicode{tfm:lmex10/integraldisplay}{222B}
\pdfglyphtounicode{tfm:lmex10/hatwide}{02C6}
\pdfglyphtounicode{tfm:lmex10/tildewide}{02DC}
\pdfglyphtounicode{tfm:lmex10/bracketleftBig}{005B}
\pdfglyphtounicode{tfm:lmex10/bracketrightBig}{005D}
\pdfglyphtounicode{tfm:lmex10/radicalbig}{221A}
\pdfglyphtounicode{tfm:lmex10/radicalBig}{221A}
\pdfglyphtounicode{tfm:lmex10/radicalBigg}{221A}
\usepackage[margin=28mm]{geometry}
\usepackage{amsmath,amssymb,amsthm,mathtools}
\usepackage{aliascnt}
\usepackage{microtype,enumitem,booktabs}
\usepackage{tikz}
\usetikzlibrary{arrows.meta,calc,positioning,fit,backgrounds,decorations.pathreplacing}
\usepackage{caption}
\usepackage{needspace}
\usepackage[colorlinks=true,linkcolor=blue!45!black,citecolor=blue!45!black,urlcolor=blue!45!black,bookmarksnumbered=true]{hyperref}
\usepackage[nameinlink,capitalise,noabbrev]{cleveref}
\pdfinfoomitdate=1
\pdftrailerid{}
\pdfsuppressptexinfo=15
\newtheorem{theorem}{Theorem}[section]
\newaliascnt{lemma}{theorem}
\newtheorem{lemma}[lemma]{Lemma}
\aliascntresetthe{lemma}
\newaliascnt{proposition}{theorem}
\newtheorem{proposition}[proposition]{Proposition}
\aliascntresetthe{proposition}
\newaliascnt{corollary}{theorem}
\newtheorem{corollary}[corollary]{Corollary}
\aliascntresetthe{corollary}
\newaliascnt{claim}{theorem}
\newtheorem{claim}[claim]{Claim}
\aliascntresetthe{claim}
\theoremstyle{definition}
\newaliascnt{definition}{theorem}
\newtheorem{definition}[definition]{Definition}
\aliascntresetthe{definition}
\theoremstyle{remark}
\newaliascnt{remark}{theorem}
\newtheorem{remark}[remark]{Remark}
\aliascntresetthe{remark}
\crefname{lemma}{Lemma}{Lemmas}
\Crefname{lemma}{Lemma}{Lemmas}
\crefname{proposition}{Proposition}{Propositions}
\Crefname{proposition}{Proposition}{Propositions}
\crefname{corollary}{Corollary}{Corollaries}
\Crefname{corollary}{Corollary}{Corollaries}
\crefname{claim}{Claim}{Claims}
\Crefname{claim}{Claim}{Claims}
\crefname{definition}{Definition}{Definitions}
\Crefname{definition}{Definition}{Definitions}
\crefname{remark}{Remark}{Remarks}
\Crefname{remark}{Remark}{Remarks}
\makeatletter
\AddToHook{env/theorem/before}{\if@nobreak\else\Needspace{4\baselineskip}\fi}
\AddToHook{env/lemma/before}{\if@nobreak\else\Needspace{4\baselineskip}\fi}
\AddToHook{env/proposition/before}{\if@nobreak\else\Needspace{4\baselineskip}\fi}
\AddToHook{env/corollary/before}{\if@nobreak\else\Needspace{4\baselineskip}\fi}
\AddToHook{env/claim/before}{\if@nobreak\else\Needspace{4\baselineskip}\fi}
\AddToHook{env/definition/before}{\if@nobreak\else\Needspace{4\baselineskip}\fi}
\AddToHook{env/remark/before}{\if@nobreak\else\Needspace{4\baselineskip}\fi}
\makeatother
\numberwithin{equation}{section}
\newcommand{\pc}{{p_c}}
\newcommand{\pu}{{p_u}}
\newcommand{\ptwo}{{p_{2\to2}}}
\newcommand{\pexp}{{p_{\exp}}}
\newcommand{\E}{\mathbb E}
\newcommand{\Prob}{\mathbb P}
\newcommand{\one}{\mathbf 1}
\newcommand{\N}{\mathbb N}
\newcommand{\Z}{\mathbb Z}
\newcommand{\R}{\mathbb R}
\newcommand{\norm}[1]{\lVert#1\rVert}
\newcommand{\abs}[1]{\lvert#1\rvert}
\newcommand{\ip}[2]{\langle#1,#2\rangle}
\DeclareMathOperator{\tr}{tr}
\DeclareMathOperator{\Adj}{Adj}
\DeclareMathOperator{\dist}{dist}
\DeclareMathOperator{\diam}{diam}
\DeclareMathOperator{\Aut}{Aut}
\DeclareMathOperator{\supp}{supp}
\DeclareMathOperator{\Nb}{Nb}
\DeclareMathOperator{\Cov}{Cov}
\DeclareMathOperator{\Var}{Var}
\setlist[enumerate]{itemsep=3pt,topsep=5pt}
\setlist[itemize]{itemsep=3pt,topsep=5pt}
\captionsetup{font=small,labelfont=bf}
\title{Nonuniqueness of percolation on\\
nonamenable quasi-transitive graphs}
\author{OpenAI}
\date{September 24, 2026}
\hypersetup{pdftitle={Nonuniqueness of percolation on nonamenable quasi-transitive graphs},pdfauthor={OpenAI},pdfsubject={The bond-percolation nonuniqueness conjecture of Benjamini and Schramm}}
\begin{document}
\maketitle
\begin{abstract}
We prove the bond-percolation nonuniqueness conjecture of Benjamini and
Schramm: every infinite connected, locally finite, nonamenable quasi-transitive
graph has a nonempty interval of parameters for which Bernoulli bond
percolation almost surely has infinitely many infinite clusters. We also prove
Hutchcroft's stronger operator-threshold conjecture, establishing
$p_c<p_{2\to2}\le p_u$.
\end{abstract}
\clearpage
\begingroup
\small
\tableofcontents
\endgroup
\clearpage
\section{Introduction}\label{sec:introduction}

Bernoulli bond percolation on an infinite graph can have two distinct
transitions. Infinite open clusters first appear at a critical parameter
$\pc$; uniqueness begins at a second threshold $\pu$. On a nonamenable
graph, the expanding geometry can keep infinite clusters apart even
after percolation has begun. Benjamini and Schramm conjectured that this
mechanism always produces a nonempty nonuniqueness phase when the graph
has finitely many vertex orbits
\cite[Conjecture~6]{BenjaminiSchramm1996}. We prove this conjecture for
Bernoulli bond percolation, together with the stronger strict inequality
for the $\ell^2$ operator threshold conjectured by Hutchcroft.
This is the equivalent $\ell^2$ formulation of his operator-threshold
conjecture \cite[Proposition~2.3 and Conjecture~6.1]{Hutchcroft2019}.

\subsection{The graph, the thresholds, and the main result}

Let $G=(V,E)$ be an infinite connected locally finite undirected graph.
It is \emph{quasi-transitive} if its automorphism group $\Aut(G)$ has
finitely many orbits on $V$. In particular its vertex degrees are
uniformly bounded. For a finite nonempty set $A\subset V$, write
\[
 \partial_V A=\{v\in V\setminus A:v\sim a\text{ for some }a\in A\},
 \qquad
 h_V(G)=\inf_{0<|A|<\infty}\frac{|\partial_V A|}{|A|}.
\]
The graph is \emph{nonamenable} if $h_V(G)>0$. For a bounded-degree
graph this is equivalent to positivity of the edge isoperimetric
constant. Loops and multiple edges may be allowed; a subdivision
reduction in \cref{pre:simple-reduction} permits the main proof to use
a simple graph.

Under $\Prob_p$, every unordered edge is independently open with
probability $p$. Write $C_x$ for the open cluster of $x$, and
$x\leftrightarrow y$ for connection by a finite open path. Paths of
length zero are allowed. If $N_\infty(p)$ is the number of infinite
open clusters, define
\begin{align*}
 \pc(G)&=\inf\{p\in[0,1]:\Prob_p(N_\infty(p)\ge1)>0\},\\
 \pu(G)&=\inf\{p\in[0,1]:\Prob_p(N_\infty(p)=1)=1\}.
\end{align*}
The two-point connection kernel is $T_p(x,y)=\Prob_p(x\leftrightarrow y)$.
We equip $V$ with counting measure and put
\[
 \ell^2(V)=\left\{f:V\to\mathbb C:\sum_{x\in V}|f(x)|^2<\infty\right\},
 \qquad (T_pf)(x)=\sum_{y\in V}T_p(x,y)f(y)
\]
initially for finitely supported $f$. Boundedness of $T_p$ means that
this action extends to a bounded operator on $\ell^2(V)$. Its threshold is
\[
 \ptwo(G)=\sup\{p\in[0,1]:T_p\text{ is bounded on }\ell^2(V)\}.
\]
All operator norms in this paper use counting measure, including
operators restricted to a vertex orbit. The weights introduced later
average the choice of root and define a finite trace; they do not
change the Hilbert-space measure.

\begin{theorem}\label{thm:main}
Let $G$ be an infinite connected locally finite quasi-transitive graph
with $h_V(G)>0$. For Bernoulli bond percolation on $G$,
\[
 \sup_{0\le p<\pc}\norm{T_p}_{2\to2}
       =\norm{T_{\pc}}_{2\to2}<\infty,
 \qquad \pc(G)<\ptwo(G)\le\pu(G).
\]
There exist $\pc<p_1<p_2<1$ such that, in the coupling obtained from
independent uniform edge labels, almost surely there are infinitely
many infinite open clusters simultaneously for every $p\in[p_1,p_2]$.
\end{theorem}

The theorem resolves the bond-percolation nonuniqueness conjecture of
Benjamini and Schramm positively, and proves Hutchcroft's stronger
operator-threshold conjecture in the same graph class. The bounds and
the interval may depend on $G$. The simultaneous conclusion is a
consequence of the strict gap, in accordance with the simultaneous
phase theorem of H\"aggstr\"om, Peres, and Schonmann
\cite[Theorem~1.3]{HaggstromPeresSchonmann1999}. We include a direct
proof for the compact interval asserted here.

\begin{corollary}[Arbitrary Cayley graphs]\label{cor:cayley}
Let $\Lambda$ be a nonamenable finitely generated group, and let $S$
be any finite symmetric generating set with the identity removed.
For Bernoulli bond percolation on the simple undirected Cayley graph
with edges $\{x,xs\}$, $x\in\Lambda$, $s\in S$, the critical connection
operator is bounded on counting-measure $\ell^2(\Lambda)$ and
\[
 \pc<\ptwo\le\pu.
\]
The simultaneous nonuniqueness interval in \cref{thm:main} exists for
every such fixed generating set. At every fixed $p\in(\pc,\pu)$,
there are almost surely infinitely many infinite open clusters.
\end{corollary}

\begin{proof}
Left translations act transitively on this locally finite graph.
The graph form of the F\o lner criterion gives positive isoperimetric
constant for every fixed finite generating set of a nonamenable group.
Apply \cref{thm:main}.
The final assertion follows from the phase theorem
\cite[Theorem~1.3]{HaggstromPeresSchonmann1999}.
\end{proof}

\subsection{Proof strategy}

For $p\in[0,1]$, write
\[
 \chi_x(p)=\E_p|C_x|=\sum_yT_p(x,y),\qquad
 \chi_{\max}(p)=\sup_x\chi_x(p),
\]
for the susceptibility and its supremum over vertices. These quantities
are finite for $p<\pc$ by sharpness (\cref{pre:sharpness}). The row sums
bound the operator norm, but that bound need not remain finite as
$p\uparrow\pc$. Our task is to control
$\norm{T_\pc}_{2\to2}$ independently of these row sums. The proof first
bounds their near-critical growth and then constructs a deformation
that reduces row sums while retaining a large operator norm.
The nonunimodular case follows from Hutchcroft's established operator
theorem. After the reductions in \cref{pre:section}, the new argument
therefore concerns a simple graph with a unimodular automorphism group.
Its critical clusters are finite by the critical-nonpercolation theorem
recalled there.

\paragraph{From local endpoint maps to critical moments.}
A local endpoint map is a randomized rule that either fails or selects
a vertex in the root's critical open cluster, with a fixed cap $M$ on
that whole cluster. The roots and selected vertices lie in one vertex
orbit. Let $u$ be its success probability, and let $K(x,y)$ be the
probability of selecting $y$ conditional on success at $x$.
If some critical moment of order
$0<\alpha<1/2$ were infinite, \cref{amp:section} would construct such
maps with $0<\norm K<1$ and
\[
 \frac{\log(1/u)+\log M}{\log(1/\norm K)}\longrightarrow0.
\]
Thus the endpoint kernel becomes dispersive much faster than the success
probability deteriorates or the cap grows.

The construction exchanges finite pieces between independent
percolation fields. The geometric input, proved in \cref{good:section},
holds outside a stretched-exponentially small event: after any vertex
deletion from a capped critical cluster, it bounds the number of bridges
separating two vertices in each remaining component that has an open
contact with the deleted set, in terms of its number of contacts.
An inverse comparison and a
joint second-moment bound in \cref{amp:section} control the multiplicity
of the exchanges and keep the output mass from concentrating on rare
fields. In \cref{mom:section}, each fixed map is run just below $\pc$.
Its small operator norm forces its iterates to send substantial endpoint
mass outside an independently sampled cluster whose law is weighted
by its size. Comparing the susceptibility in that deleted
domain with the full susceptibility gives a differential inequality.
A finite-exploration comparison with critical percolation then
contradicts the assumed divergent moment. This proves
\[
 \sup_x\E_\pc|C_x|^\alpha<\infty\quad(0<\alpha<1/2).
\]
A separate decision-tree argument in the same section converts these
moments into
\[
 \chi_{\max}(p)\le C_\eta(\pc-p)^{-1-\eta}
 \qquad\text{for every }\eta>0.
\]

\paragraph{A subpower triangle bound.}
The next quantity is the triangle diagram
\[
 \nabla_p(x)=\sum_{y,z}T_p(x,y)T_p(y,z)T_p(z,x).
\]
\Cref{trace:section} compares two explorations of a cluster, one in the
full graph and one after another cluster has been removed. The
difference of their centered edge scores---sums of tested edge bits
after subtracting their means---counts the closed contacts between
the clusters. Concentration of the scores and the fractional
moments above bound those contacts. A differential inequality for the
type average of the triangle diagram then yields
\[
 \sup_x\nabla_p(x)\le C_\eta(\pc-p)^{-\eta}
 \qquad\text{for every }\eta>0.
\]
The supremum and type-average forms are equivalent up to fixed constants
because there are finitely many vertex types. This subpower bound is the
remaining input to the deformation argument.

\paragraph{Corridor deformation and the spectral contradiction.}
A corridor is a simple open path whose internal vertices have no other
open incident edges. For fixed $p<\pc$, \cref{def:section} penalizes a
connection for each specified corridor all of whose edges are pivotal
for it. Raw paths of length $k$ from a symmetric lazy walk specify these
corridors by loop erasure; the penalty weights cancel the probabilities
of forcing their open edges and closed boundaries.

Nonamenability gives the walk an $\ell^2$ norm $\rho<1$. After deformation
time $t$, the norm of the connection kernel loses at most
$2t\rho^k\norm{T_p}^2$. Up to an explicit time-dependent factor, the
rate of decrease of its averaged row sum is bounded below by an average
of products of the deformed susceptibilities, or weighted row sums, at the two
corridor endpoints in the graph with the corridor interior deleted,
apart from an error
bounded by a constant times
$\chi_{\max}(p)^2\rho^k\sup_x\nabla_p(x)$.
\Cref{path:section} finds a positive proportion of raw paths for which
both endpoint row sums remain substantial. Geodesic segments
at the two ends provide escapes from the raw path, and random-walk
smoothing makes the effect of deleting the corridor small at the
escape endpoints.

If $T_\pc$ were unbounded, sprinkling would imply
$\norm{T_p}\ge c(\pc-p)^{-1}$. The preceding susceptibility bound
would make $\chi_{\max}(p)/\norm{T_p}$ subpower as well. In
\cref{con:section}, these estimates permit a length
$k=o(\log(1/(\pc-p)))$ for which the norm stays large while the
averaged row sum decreases quadratically in its own size. The row sum
would become too small to bound the surviving operator norm. This
contradiction proves critical boundedness. Sprinkling then gives
$\pc<\ptwo$, the uniqueness comparison gives $\ptwo\le\pu$, and a
finite-modification argument supplies the simultaneous interval in
\cref{thm:main}.

\subsection{Critical behavior and exponential connectivity}

The operator conclusion gives more than the separation of the two
percolation thresholds. It also determines the powers governing the
size and radius of critical clusters, and ensures that two-point
connections still decay exponentially in part of the supercritical
phase. We state these consequences using the following observables.
The triangle diagram $\nabla_p(v)$ and susceptibility $\chi_v(p)$
were defined above; either may be infinite. For $p\in[0,1]$ and
$v\in V$, write
\[
 \theta_v(p)=\Prob_p(|C_v|=\infty),
\]
and let $G[p]$ be the open spanning subgraph.
The extrinsic and intrinsic radii of the cluster of $v$ are
\[
 \operatorname{rad}_{\mathrm{ext}}(C_v)
    =\sup_{x\in C_v}\dist_G(v,x),\qquad
 \operatorname{rad}_{\mathrm{int}}(C_v)
    =\sup_{x\in C_v}\dist_{G[p]}(v,x),
\]
with values in $[0,\infty]$. Thus the intrinsic radius measures
distance along open paths, whereas the extrinsic radius uses the
ambient graph.

Say that connection probabilities at $p$ decay exponentially if
there are constants $0<A<\infty$ and $a>0$ such that
$T_p(x,y)\le A e^{-a\dist_G(x,y)}$ for every $x,y\in V$, and define
\[
 \pexp(G)=\sup\{p\in[0,1]:
       \text{connection probabilities at $p$ decay exponentially}\}.
\]
This is the exponential-connectivity threshold of
\cite[Section~2]{Hutchcroft2020L2}. We write $f\asymp g$ in a stated
limit when $c g\le f\le Cg$ sufficiently near that limit for some
constants $0<c\le C<\infty$.

\Needspace{25\baselineskip}
\begin{corollary}[Critical laws and exponential connectivity]
\label{cor:critical-laws}
Let $G$ be an infinite connected locally finite quasi-transitive graph
with $h_V(G)>0$. For Bernoulli bond percolation on $G$,
\begin{equation}\label{eq:critical-triangle}
 \nabla_{\pc}(v)\le\norm{T_{\pc}}_{2\to2}^{3}<\infty
 \qquad\text{for every }v\in V.
\end{equation}
For every $v\in V$,
\begin{align}
 \chi_v(p)&\asymp(\pc-p)^{-1} &&(p\uparrow\pc), \label{eq:meanfield-chi}\\
 \theta_v(p)&\asymp p-\pc &&(p\downarrow\pc), \label{eq:meanfield-theta}\\
 \Prob_{\pc}(|C_v|\ge n)&\asymp n^{-1/2} &&(n\to\infty), \label{eq:meanfield-size}\\
 \Prob_{\pc}\bigl(\operatorname{rad}_{\mathrm{int}}(C_v)\ge n\bigr)
    &\asymp n^{-1} &&(n\to\infty), \label{eq:meanfield-intrinsic}\\
 \Prob_{\pc}\bigl(\operatorname{rad}_{\mathrm{ext}}(C_v)\ge n\bigr)
    &\asymp n^{-1} &&(n\to\infty). \label{eq:meanfield-extrinsic}
\end{align}
Here $n$ tends through positive integers. The comparison constants
and the one-sided neighborhoods of $\pc$ may depend on $G$ and $v$.
Moreover,
\begin{equation}\label{eq:exponential-threshold}
 \pc(G)<\ptwo(G)\le\pexp(G).
\end{equation}
For each fixed $p\in[0,\ptwo(G))$, there are constants
$0<A_{G,p}<\infty$ and $a_{G,p}>0$ such that
\begin{equation}\label{eq:exponential-connectivity}
 T_p(x,y)\le A_{G,p}e^{-a_{G,p}\dist_G(x,y)}
 \qquad\text{for every }x,y\in V.
\end{equation}
\end{corollary}

The triangle-condition assertion and the strict inequality
$\pc<\pexp$ verify Conjectures~1.2 and~2.1 of
\cite{Hutchcroft2020L2}. In particular, exponential connectivity
persists at $\pc$ and throughout the nonempty interval
$(\pc,\ptwo)$, even though infinite clusters already exist there.
The power laws are applications of earlier conditional theorems:
the new input is that the operator condition holds on every graph in
the stated class. We give the deduction, including the graph
conventions and the distinction between intrinsic and extrinsic
radius, in \cref{cons:section}.

\subsection{Historical context and the methods used here}

The distinction between existence of an infinite cluster and its
uniqueness is central to the program initiated by Benjamini and Schramm
\cite{BenjaminiSchramm1996}. Their formulation uses almost transitive
graphs, meaning the quasi-transitive graphs considered here.
Their standing model is site percolation, and they explicitly extend
their nonplanar questions to bond percolation on page~72. A regular
tree of degree at least three supplies the basic example of separated
thresholds. The difficulty is to extract an equally robust separation
from nonamenability alone, without a tree structure or an additional
geometric hypothesis.

Several earlier results identify additional geometry that guarantees
nonuniqueness. Benjamini and Schramm proved $0<\pc<\pu<1$ for
one-ended nonamenable transitive planar graphs, for both bond and site
percolation \cite[Theorem~1.1]{BenjaminiSchramm2001}. Their result
followed Lalley's work on site percolation on canonical graphs of
cocompact Fuchsian groups \cite{Lalley1998}. Schonmann established
multiple phase transitions in a quantitatively highly nonamenable
regime \cite{Schonmann2001}. Nachmias and Peres proved nonuniqueness
and triangle estimates for nonamenable transitive graphs of sufficiently large
girth relative to their spectral radius \cite{NachmiasPeres2012}.
These results explain the usefulness of geometric separation and
random-walk decay, while imposing conditions beyond nonamenability.

A different approach uses harmonic functions of finite Dirichlet energy.
Gaboriau proved $\pc<\pu$ for unimodular transitive graphs admitting a
nonconstant harmonic Dirichlet function, using $\ell^2$ Betti numbers,
and extended the argument to quasi-transitive graphs
\cite[Theorem~0.6 and Section~5]{Gaboriau2005}. This supplies a further
class in which nonuniqueness follows from a hypothesis stronger than
nonamenability alone.

For Cayley graphs, the quantifier over generating sets is significant.
Pak and Smirnova-Nagnibeda proved nonuniqueness for a suitable generating
multiset \cite{PakSmirnovaNagnibeda2000}. Thom proved that a suitable ordinary
finite symmetric generating set exists \cite[Corollary~8]{Thom2015}. These results
permit a choice of generators; \cref{cor:cayley} treats an arbitrary
prescribed generating set. The quasi-transitive theorem also permits
several vertex orbits and nontrivial vertex stabilizers, for which group
coordinates are not available.

Critical nonpercolation is a separate established input.
Benjamini, Lyons, Peres, and Schramm proved that critical percolation on
every nonamenable Cayley graph has no infinite clusters, and explained
the extension to unimodular quasi-transitive graphs
\cite[Theorem~1.1 and Section~2]{BLPS1999}. Hutchcroft established this conclusion on
every connected locally finite quasi-transitive graph of exponential
growth \cite[Theorem~1.2]{Hutchcroft2016}. Positive isoperimetry implies
exponential growth, so the latter theorem gives the finite critical
clusters on which our local constructions operate.

The behavior of infinite clusters in a common coupling has its own
history. H\"aggstr\"om and Peres proved uniqueness monotonicity and
simultaneous phase statements under unimodularity
\cite{HaggstromPeres1999}. Schonmann established stability and
uniqueness monotonicity for all quasi-transitive graphs
\cite{Schonmann1999}. H\"aggstr\"om, Peres, and Schonmann then proved
simultaneous phase counts away from the two threshold values in that
generality \cite[Theorem~1.3]{HaggstromPeresSchonmann1999}. These
results identify the probabilistic consequences of a strict threshold
gap; establishing that gap is the task here.

Hutchcroft introduced and developed the operator-threshold framework
used here, proved $\pc<\ptwo$ for nonamenable Gromov hyperbolic
quasi-transitive graphs, and proved it for graphs admitting a
quasi-transitive nonunimodular automorphism group
\cite[Section~2 and Theorems~2.1 and~2.9]{Hutchcroft2019}.
The nonunimodular result builds on his tilted approach to
nonuniqueness and mean-field criticality
\cite[Theorem~1.2]{Hutchcroft2020Nonunimodular}; the stronger operator
form used here is Theorem~2.9 of \cite{Hutchcroft2019}.
Choi and Seo subsequently obtained nonuniqueness and the triangle
condition for every Cayley graph of a finitely generated acylindrically
hyperbolic group by extending Hutchcroft's geometric criteria
\cite[Theorems~A and~B]{ChoiSeo2025}.
The nonunimodular theorem is used directly in our reduction. The new
argument therefore concerns a unimodular action. Its averaging over
vertex types uses the inverse stabilizer-volume weights from the
mass-transport theory of Benjamini, Lyons, Peres, and Schramm
\cite[Corollary~3.5]{BLPS1999Invariant}. The corresponding root law
and operator trace fit the general framework of Aldous and Lyons
\cite[Theorem~3.1 and Section~5]{AldousLyons2007}. We prove the
specific transport and trace identities needed below.

Hutchcroft developed the critical and supercritical consequences of
the operator condition in \cite{Hutchcroft2020L2}. His
Theorem~2.2 compares the operator and exponential-connectivity
thresholds, and his Theorem~3.1 gives the critical extrinsic-radius
law under the strict operator gap. The critical intrinsic-radius
upper bound has its origin in the lattice theorem of Kozma and
Nachmias \cite[Theorem~1.2(ii)]{KozmaNachmias2009}.
Hutchcroft's later finite-cluster theorem supplies both radius
estimates in the quasi-transitive setting under the operator
condition \cite[Theorem~1.2]{Hutchcroft2022FiniteClusters}.
These are the conditional geometric results applied in
\cref{cons:section}.
The upper comparison $\ptwo\le\pu$ can be an equality.
Hutchcroft and Pan proved $\ptwo=\pu$ for the Cartesian
product of a regular tree of degree at least three with an infinite
amenable Cayley graph
\cite[Theorem~1.2]{HutchcroftPan2024Dimension}. Thus the strict lower
comparison established here is compatible with a sharp upper endpoint.

Our probabilistic estimates use Harris's association inequality and
the van den Berg--Kesten inequality for disjoint open witnesses
\cite{Harris1960,VanDenBergKesten1985}. Subcritical susceptibility
finiteness is obtained by adapting the finite-set sharpness argument
of Duminil-Copin and Tassion \cite[Section~1]{DuminilCopinTassion2016}
to finitely many vertex types. Sharpness originates in the independent
work of Aizenman--Barsky and Menshikov
\cite{AizenmanBarsky1987,Menshikov1986}; Antunovi\'c and Veseli\'c
proved the full quasi-transitive susceptibility and exponential-tail
statements \cite{AntunovicVeselic2008}. The subsequent conversion of critical
cluster-size estimates into near-critical susceptibility estimates
uses independent vertex marks and a decision-tree covariance
inequality of O'Donnell, Saks, Schramm, and Servedio
\cite[Section~3.3, Equation~(6)]{ODonnellSaksSchrammServedio2005}.
This is the same critical-tail-to-susceptibility principle developed
by Hutchcroft \cite[Theorem~1.1]{Hutchcroft2020Exponents}; we give
the needed proof uniformly over vertex types and induced domains.
The sprinkling comparison that finally opens an interval above
$\pc$ is the connection-operator expansion of
\cite[Lemma~2.4]{Hutchcroft2019}, also proved here.

The contact estimates also have a close methodological antecedent in
Hutchcroft's two-ghost inequalities and their critical-tail refinements
\cite{Hutchcroft2020Locality,Hutchcroft2021PowerLaw}. Those arguments
use independent marks, cluster exploration, and centered edge scores,
following the fluctuation method of Aizenman, Kesten, and Newman
\cite{AizenmanKestenNewman1987}. Here we compare a full-cluster
exploration with an exploration in a deleted domain to control the
number of closed contacts between two clusters. The triangle diagram
and positive connection matrix belong to the tree-graph framework of
Aizenman and Newman \cite{AizenmanNewman1984}; the triangle-condition
work of Barsky--Aizenman and the subsequent quantitative development
by Hutchcroft provide further context
\cite{BarskyAizenman1991,Hutchcroft2022Triangle}. Our required subpower
trace bound and its use in the corridor comparison are proved below.

Finite modifications that join separated clusters, create pivotal edges,
and control the multiplicity of their preimages already appear in Kozma's
work on products of trees \cite[Section~1.4 and Lemmas~7--11]{Kozma2011}.
Our local endpoint maps exchange finite pieces between independent
percolation fields; the deletion-uniform bridge bounds and paired inverse
estimate needed for these exchanges are proved here. Together with the
deformation that penalizes forced pivotal corridors, these maps let us
compare the spreading of a connection kernel in $\ell^2$ with the total
mass of its rows. Both constructions are formulated through covariant
kernels on vertex orbits.

\section{Preliminaries and the unimodular reduction}
\label{pre:section}

Throughout the proof, all operator norms refer to counting-measure
$\ell^2(V)$. A kernel is called invariant if its value at $(gx,gy)$ equals
its value at $(x,y)$ for every automorphism $g$. For a vertex set $D$, the
same letter denotes the induced subgraph on $D$. Write $C_x^D$ for the
open cluster of $x$ in $D$, and put
\[
 C_x=C_x^V,\qquad s_x=|C_x|,\qquad
 T_p^D(x,y)=\Prob_p(x\leftrightarrow y\text{ in }D),\qquad T_p=T_p^V.
\]
Quantities requiring a vertex outside their domain are defined to be zero.
We use $\partial_E A$ for the set of edges with exactly one endpoint in
$A$, and $\Nb[A]$ for $A$ together with all its neighbors.

\subsection{Graph reductions and the two established inputs}
\label{pre:reductions}

\begin{lemma}[Reduction to simple graphs]\label{pre:simple-reduction}
It suffices to prove the operator gap and the infinite-cluster conclusion
for simple graphs satisfying the hypotheses of \cref{thm:main}.
\end{lemma}
\begin{proof}
Loops never affect connectivity and may be removed. Replace each remaining
edge of $G$, including each copy of a multiple edge, by a two-edge path
through its own new vertex, obtaining a simple graph $\widetilde G$.
Quasi-transitivity and local finiteness give a finite degree bound $d_0$
and finitely many edge orbits in $G$. The lifted automorphisms consequently
have finitely many vertex orbits in $\widetilde G$.

To verify nonamenability, let a finite set $A$ in $\widetilde G$ contain
$n$ old vertices and $m$ new vertices. Each distinct boundary neighbor in
$G$ of its old-vertex set supplies a distinct boundary vertex of $A$:
along a selected subdivided edge, use the new vertex if it is outside
$A$, and otherwise use the old boundary neighbor. Hence
$|\partial_V A|\ge h_V(G)n$. At most $d_0n$ selected new vertices have
an endpoint among the selected old vertices. If $m>2d_0n$, more than
$m/2$ selected new vertices have both endpoints outside $A$; bounded
degree then gives $|\partial_V A|\ge m/(2d_0)$. The first estimate when
$m\le2d_0n$ and the second otherwise give a positive lower bound for
$|\partial_V A|/(n+m)$, independent of $A$.

In percolation of parameter $t$ on $\widetilde G$, declare an original
edge open precisely when both its replacement edges are open. These
independent declarations have parameter $t^2$. Infinite clusters
correspond bijectively: a cluster with only finitely many old vertices
is finite by local finiteness. In particular
\[
 \pc(\widetilde G)^2=\pc(G),\qquad
 \pu(\widetilde G)^2=\pu(G).
\]
Moreover, $T_{t^2}^G$ is the old-vertex compression of
$T_t^{\widetilde G}$. Boundedness of the latter implies boundedness of
the former. Thus a strict operator gap, and a coupled interval of
infinitely many infinite clusters, both transfer to $G$.
\end{proof}

We henceforth assume that $G$ is simple and fix an integer $d$ bounding
its degrees. Write $\Adj(x,y)=\one_{\{x\sim y\}}$ for its adjacency kernel.
Positive vertex isoperimetry supplies a constant $\iota>0$
such that $|\partial_E A|\ge\iota|A|$ for every finite nonempty $A$.

\begin{lemma}[Critical finiteness]\label{pre:critical-finiteness}
We have $0<\pc<1$. At $\pc$, all open clusters are finite almost surely,
simultaneously at every vertex and in every induced subgraph under
restriction of the full configuration.
\end{lemma}
\begin{proof}
There are at most $d^n$ length-$n$ paths from a fixed vertex, so no
infinite cluster exists when $pd<1$. For the opposite bound, the number
of connected vertex sets of size $n$ containing a fixed vertex is at
most $d^{2n}$: choose a deterministic spanning tree of each set and
encode it by a traversal of length $2n-2$. If such a set is the root
cluster, all its boundary edges are closed. Consequently
\[
 \Prob_p(|C_x|<\infty)
 \le \sum_{n\ge1}d^{2n}(1-p)^{\iota n}<1
\]
when $p$ is sufficiently close to $1$.

For every graph-distance ball,
\[
 |B(x,r+1)|\ge(1+h_V(G))|B(x,r)|,
\]
so $G$ has exponential growth. Hutchcroft's critical-finiteness theorem
for connected locally finite quasi-transitive graphs of exponential
growth \cite[Theorem~1.2]{Hutchcroft2016} therefore applies to bond
percolation on $G$. Countability gives simultaneous finiteness at all
vertices. Every cluster in a restricted configuration is contained in
a full-graph cluster, proving the last assertion.
\end{proof}

\subsection{Product inequalities and operator bounds}
\label{pre:product-tools}

The first two assertions below are Harris's association inequality
\cite{Harris1960} and the van den Berg--Kesten (BK) inequality
\cite{VanDenBergKesten1985}. We include the product-measure proofs.

\begin{lemma}[Association, disjoint witnesses, and the tree bound]
\label{pre:probability-tools}
Increasing nonnegative functions of independent percolation bits have
nonnegative covariance. If increasing events have disjoint finite open
witnesses, their joint disjoint-occurrence probability is at most the
product of their probabilities. In particular, for all $x,y,z$,
\begin{equation}\label{eq:s01}
 \Prob_p(x\leftrightarrow y,\ x\leftrightarrow z)
 \le\sum_w T_p(x,w)T_p(w,y)T_p(w,z).
\end{equation}
The same assertions hold in every induced domain.
\end{lemma}
\begin{proof}
For association in a finite product, condition on the last bit and use
induction. The two conditional means are increasing functions of that
bit, whose covariance is nonnegative. Truncation and conditional
expectations on increasing finite coordinate sets extend this argument
to the countable product and to nonnegative functions.

Here is a finite-product proof of the disjoint-witness inequality.
For two increasing events $A,B$, denote disjoint occurrence by
$A\mathbin{\Box}B$. Suppose first that $p\le1/2$, and let $A_0,A_1$
and $B_0,B_1$ be the sections at the last bit. Write
$D_{ij}=A_i\mathbin{\Box}B_j$. The zero-section of
$A\mathbin{\Box}B$ is contained in $D_{00}$; its one-section is
contained in $D_{10}\cup D_{01}$. Also
$D_{00}\subseteq D_{10}\cap D_{01}$. Thus
\[
 \Prob(A\mathbin{\Box}B)
 \le(1-2p)\Prob(D_{00})+p\Prob(D_{10})+p\Prob(D_{01}).
\]
Apply induction, putting $a_i=\Prob(A_i)$ and $b_i=\Prob(B_i)$.
The difference between $\Prob(A)\Prob(B)$ and the resulting upper
bound is $p^2(a_1-a_0)(b_1-b_0)\ge0$.
For $1/2<p<1$, realize each bit as the logical OR of $m$ independent
bits of parameter $q=1-(1-p)^{1/m}\le1/2$. Disjoint open witnesses
lift to disjoint open witnesses in the enlarged product. The cases
$p=0,1$ are immediate. Iteration proves the assertion for any finite
number of events, and exhaustion proves the finite-witness assertion
on $G$.

On $\{x\leftrightarrow y,\ x\leftrightarrow z\}$, take a finite open
tree joining the three vertices. Its branch vertex $w$ supplies three
edge-disjoint paths to $x,y,z$, with zero-length paths permitted.
A union bound over $w$, followed by the disjoint-witness inequality,
gives \eqref{eq:s01}.
\end{proof}

The tree bound is the three-terminal instance of the tree-diagram
method of Aizenman and Newman \cite{AizenmanNewman1984}.
We repeatedly use two elementary kernel facts. If $|A(x,y)|\le B(x,y)$
and $B$ is a nonnegative bounded kernel, then $\norm A\le\norm B$:
test the bilinear form against absolute values of finitely supported
vectors. If all row sums and column sums of $B$ are at most $b$, then
$\norm B\le b$, by Cauchy--Schwarz with the weights $B(x,y)$.

\begin{lemma}[The operator threshold precedes uniqueness]
\label{pre:threshold-order}
For every quasi-transitive graph under consideration,
$\ptwo\le\pu$.
\end{lemma}
\begin{proof}
Suppose that there is a unique infinite cluster almost surely at
parameter $p$. Some vertex has positive probability of belonging to
it, by countability. Association with the event that a fixed finite
connecting path is open makes this probability positive at every
vertex. There are finitely many vertex types, so
$\inf_x\Prob_p(|C_x|=\infty)=\theta>0$. Association and uniqueness give
$T_p(x,y)\ge\theta^2$ for every pair $x,y$. For finite $F\subset V$,
\[
 \ip{\one_F}{T_p\one_F}\ge\theta^2|F|^2,
\]
which rules out boundedness on $\ell^2(V)$. Boundedness parameters
form a downward interval by entrywise domination. Thus $\ptwo$ is
at most every parameter of almost-sure uniqueness, proving the claim
under the stated supremum and infimum definitions.
\end{proof}

\begin{proposition}[Reduction to a unimodular action]
\label{pre:unimodular-reduction}
If $\Aut(G)$ is nonunimodular, then $\pc<\ptwo\le\pu$ and
$T_\pc$ is bounded.
For the proof of the operator gap it therefore remains to treat the
case in which $\Gamma=\Aut(G)$ is unimodular.
\end{proposition}
\begin{proof}
Give $\Aut(G)$ the topology of pointwise convergence on vertices.
Vertex stabilizers are open and compact: restrictions to successive
finite balls give the usual diagonal compactness argument. Thus
$\Aut(G)$ is a locally compact group. If it is nonunimodular, its
quasi-transitive action satisfies the hypotheses of
\cite[Theorem~2.9]{Hutchcroft2019}, which gives
$\pc<p_{q\to q}$ for every $1<q<\infty$. Take $q=2$ and apply
\cref{pre:threshold-order}. A bounded kernel at any parameter strictly
between $\pc$ and $\ptwo$ dominates $T_\pc$, proving its boundedness.
This cited theorem requires a
quasi-transitive nonunimodular automorphism subgroup; hyperbolicity
is not among its hypotheses.
\end{proof}

The remaining assertion, proved in \cref{con:section}, is the principal
new estimate.

\begin{theorem}[Critical boundedness for a unimodular action]
\label{thm:critical-operator}
Let $G$ be an infinite connected locally finite simple quasi-transitive
graph with $h_V(G)>0$. If $\Aut(G)$ is unimodular, then
$\norm{T_{\pc}}_{2\to2}<\infty$.
\end{theorem}

\subsection{Averaging over types and a finite trace}
\label{pre:types}

The inverse stabilizer-volume weights below give the quasi-transitive
mass-transport principle of \cite[Corollary~3.5]{BLPS1999Invariant}; see
also \cite[Theorem~3.1]{AldousLyons2007} and
\cite[Corollary~8.11 and Section~8.2]{LyonsPeres2016}.

For the rest of the main argument assume that $\Gamma$ is unimodular.
Let $O_1,\ldots,O_{h_0}$ be its vertex orbits, with representatives
$o_1,\ldots,o_{h_0}$. Fix a Haar measure on $\Gamma$, and set
\[
 \sigma_i=
 \frac{\operatorname{vol}(\Gamma_{o_i})^{-1}}
 {\sum_{j=1}^{h_0}\operatorname{vol}(\Gamma_{o_j})^{-1}},
 \qquad [F(o)]=\sum_{i=1}^{h_0}\sigma_iF(o_i),
 \qquad c_\sigma=(\min_i\sigma_i)^{-1}.
\]
Brackets always mean this finite type average, including when $F$
already contains a probabilistic expectation. These weights determine a
law on the root types, not a measure replacing counting measure on $V$.
In particular all adjoints, positivity assertions, and operator norms
below are those of the previously defined Hilbert space $\ell^2(V)$.

\begin{lemma}[Weighted mass transport]\label{pre:mass-transport}
For every invariant nonnegative function $m$ on ordered vertex pairs,
\begin{equation}\label{eq:s02}
 \left[\sum_y m(o,y)\right]
 =\left[\sum_x m(x,o)\right].
\end{equation}
The identity also applies to the expectation of a covariant random
mass and to absolutely summable signed or complex masses.
In particular, an invariant stochastic kernel on a single orbit has
column sums one, whether or not it is symmetric.
\end{lemma}
\begin{proof}
Decompose the ordered pairs into $\Gamma$-orbits. For one such orbit,
represented by $(x,y)$ of types $i,j$, its number of pairs with first
coordinate $x$ is
\[
 \frac{\operatorname{vol}(\Gamma_x)}
 {\operatorname{vol}(\Gamma_x\cap\Gamma_y)},
\]
and its number with second coordinate $y$ is the same expression with
$\Gamma_y$ in the numerator. Unimodularity makes the volumes of
conjugate compact subgroups equal. Multiplication by $\sigma_i$ and
$\sigma_j$ therefore gives the same contribution. Summing proves
\eqref{eq:s02}; Tonelli's theorem handles nonnegative infinite sums.
Take expectations for random masses, and use absolute convergence for
signed or complex ones. Supporting the mass on one orbit cancels its
positive type weight, giving the final assertion.
\end{proof}

The diagonal average gives a normalized trace even though $V$ is
infinite. This is the invariant-operator trace discussed in
\cite[Section~5]{AldousLyons2007}; we verify its precise form here.

\begin{lemma}[The invariant trace]\label{pre:trace}
For a bounded invariant operator $S$, define
$\tr S=[S(o,o)]$. This is a positive normalized trace:
$\tr I=1$ and $\tr(AB)=\tr(BA)$ for bounded invariant $A,B$.
If $A$ is positive semidefinite, invariant, and bounded, and $B=B^*$
is invariant and bounded,
then
\[
 |\tr(BA)|\le\norm B\,\tr A,
 \qquad A(x,x)\le c_\sigma\tr A.
\]
\end{lemma}
\begin{proof}
Positivity and normalization follow from the diagonal definition.
The diagonal product expansion is absolutely summable, since
\[
 \sum_y|A(x,y)B(y,x)|
 \le\norm{A^*\one_{\{x\}}}\,
       \norm{B\one_{\{x\}}}\le\norm A\norm B.
\]
Apply \eqref{eq:s02} to these products to obtain cyclicity.
The positive square root $A^{1/2}$ is invariant, because $A$ commutes
with the unitary action of every automorphism and functional calculus
preserves that commutation. The inequalities
\[
 -\norm B A\le A^{1/2}BA^{1/2}\le\norm B A
\]
and cyclicity prove the trace estimate. Finally the nonnegative
diagonal is constant on each type, so its value on type $i$ is at
most $\sigma_i^{-1}\tr A$.
\end{proof}

\subsection{A symmetric walk with a spectral gap}
\label{pre:walk-subsection}

We use a symmetric lazy version of the walk so that counting measure
is reversible even when degrees differ between vertex types. The
spectral estimate below is the discrete Cheeger argument; see
\cite{Dodziuk1984} for its isoperimetric background. Define
\[
 P(x,y)=\frac1{2d}\quad(x\sim y),\qquad
 P(x,x)=1-\frac{\deg(x)}{2d},
\]
with all other entries zero.

\begin{lemma}[Walk and type chain]\label{pre:walk}
The kernel $P$ is positive semidefinite and
\begin{equation}\label{eq:s03}
 \rho:=\norm P<1.
\end{equation}
Its projection onto the vertex types is an irreducible lazy finite
Markov chain. There exist $c_*>0$ and $l_0<\infty$ such that
\[
 \sum_{y\in O_j}P^l(x,y)\ge c_*
 \quad\text{for every }x\in V,\ j\le h_0,\ l\ge l_0.
\]
\end{lemma}
\begin{proof}
The kernel $2P-I$ is symmetric and stochastic, hence a contraction;
therefore $P$ is positive semidefinite. For a finitely supported
nonnegative $f$, the level-set identity and edge isoperimetry give
\[
 \sum_{\{x,y\}\in E}|f(x)^2-f(y)^2|
 =\int_0^\infty|\partial_E\{f^2>t\}|\,dt
 \ge\iota\norm f^2.
\]
Cauchy--Schwarz and
$\sum_{\{x,y\}\in E}(f(x)+f(y))^2\le2d\norm f^2$ imply
\[
 \ip f{(I-P)f}
 =\frac1{2d}\sum_{\{x,y\}\in E}(f(x)-f(y))^2
 \ge\frac{\iota^2}{4d^2}\norm f^2.
\]
Since $P$ has nonnegative entries,
$|\ip f{Pf}|\le\ip{|f|}{P|f|}$, so the same norm bound extends to
arbitrary finitely supported vectors and then to $\ell^2(V)$.

The transition probabilities between types are well defined by
invariance. Connectivity gives irreducibility, and every holding
probability is at least $1/2$. By following a connecting type path
and padding it with holds, a sufficiently large common power of this
finite transition matrix has every entry positive. Its minimum entry
is a lower bound for every later power, because left multiplication
by a stochastic matrix takes convex combinations of its rows.
\end{proof}

\subsection{Subcritical susceptibility, differentiation, and sprinkling}
\label{pre:subcritical}

Set
\[
 \chi_x(p)=\E_p|C_x|,\qquad
 \chi(p)=[\chi_o(p)],\qquad
 \chi_{\max}(p)=\sup_x\chi_x(p).
\]
The parameter will be suppressed when unambiguous.

\begin{lemma}[Subcritical finiteness and type comparison]
\label{pre:sharpness}
For $p<\pc$, $\chi_{\max}(p)<\infty$ and
\[
 \sup_x\E_p s_x^2\le\chi_{\max}(p)^3,
 \qquad \chi_{\max}(p)\le c_\sigma\chi(p).
\]
There is $c>0$, depending only on $G$, such that
$\chi_x(p)\ge c\chi(p)$ for every $x$ and
$\pc/2\le p<\pc$. The operator $T_p$ is symmetric, positive
semidefinite, and bounded with $\norm{T_p}\le\chi_{\max}(p)$.
\end{lemma}
\begin{proof}
Susceptibility finiteness on quasi-transitive graphs is established in
\cite[Theorem~2]{AntunovicVeselic2008}, following the sharpness work of
Aizenman--Barsky and Menshikov \cite{AizenmanBarsky1987,Menshikov1986}.
We give a proof by adapting the finite-set
criterion of Duminil-Copin and Tassion
\cite[Section~1]{DuminilCopinTassion2016} to the finitely many vertex
types.
For a finite set $S\ni x$, put
\[
 \varphi_p(x,S)=p\sum_{\substack{a\sim b\\a\in S, b\notin S}}
                 \Prob_p(x\leftrightarrow a\text{ in }S),
\]
and let $p_*$ be the supremum of parameters for which every vertex has
such an $S$ with $\varphi_p(x,S)<1$. This set of parameters is
downward closed. If $p<p_*$, choose a qualifying larger parameter;
finitely many orbit representatives then supply translated test sets
with a common size bound $L$ and costs at $p$ bounded by $a<1$.

For a finite induced domain $\Lambda$, let $H_\Lambda$ be its largest
susceptibility. A simple open path from $x\in\Lambda$ to a vertex
outside its test set $S$ has a first exit edge from $S$. Its prefix
inside $S$, that edge, and its remaining leg have disjoint witnesses.
The disjoint-witness inequality gives
\[
 \E_p|C_x^\Lambda|\le L+\varphi_p(x,S)H_\Lambda
 \le L+aH_\Lambda.
\]
This remains valid if $S$ is not contained in $\Lambda$, since the
prefix inside $S\cap\Lambda$ is dominated by the connection in $S$.
Maximizing and exhausting $G$ gives $\chi_{\max}(p)\le L/(1-a)$.
Thus $\pc\ge p_*$.

For the reverse inequality, choose $p_*<r<p<1$. There is a vertex $x$
such that $\varphi_r(x,S)\ge1$ for every finite $S\ni x$. Let
$f_\Lambda(u)$ be the probability that $x$ connects outside a finite
set $\Lambda\ni x$. This is a finite-edge event. Finite-product
differentiation expresses its derivative as the expected number of
closed pivotal edges divided by $1-u$.
Explore from outside $\Lambda$, and let $S$ be the vertices of
$\Lambda$ not connected to its exterior. On disconnection $x\in S$.
Conditional on this exploration, the edges internal to $S$ remain
independent with parameter $u$, and its crossing edges are closed.
A crossing edge is pivotal precisely when its inner endpoint connects
to $x$ in $S$. Hence, for $r\le u\le p$,
\[
 f_\Lambda'(u)
 =\frac1{u(1-u)}
   \E_u\bigl[\one_{\{x\in S\}}\varphi_u(x,S)\bigr]
 \ge\frac{1-f_\Lambda(u)}{u(1-u)}.
\]
Integration gives
$f_\Lambda(p)\ge1-\exp(-\int_r^pdu/(u(1-u)))>0$, uniformly in
$\Lambda$. Exhaustion yields $\pc\le p$; letting $p\downarrow p_*$
proves $p_*=\pc$.

Sum \eqref{eq:s01} over $y,z$ to obtain the second moment bound.
The upper type comparison follows directly from the weights.
There is a common finite bound on the length needed to reach a vertex
of any prescribed type from any starting vertex: choose paths between
representatives and translate them. Association with such a path of
length $l$ gives $\chi_x(p)\ge p^l\chi_y(p)$, proving the lower
comparison on the stated parameter interval.
Symmetry of $T_p$ and the row/column sum bound prove boundedness.
For every finitely supported complex $f$,
\[
 \sum_{x,y}\overline{f(x)}T_p(x,y)f(y)
 =\E_p\sum_{C}\left|\sum_{x\in C}f(x)\right|^2\ge0,
\]
where the sum is over open clusters. This proves positive
semidefiniteness.
\end{proof}

\begin{lemma}[Russo identities in operator norm]\label{pre:russo}
The map $p\mapsto T_p$ is locally absolutely continuous in operator
norm on $(0,\pc)$, and $\chi$ is locally absolutely continuous there.
For almost every $p\in(0,\pc)$,
\begin{equation}\label{eq:s04}
 \begin{aligned}
 (1-p)T'_p(x,y)
 &=\sum_a\sum_{b\sim a}
   \Prob_p(x\in C_a,\ y\in C_b,\ C_a\ne C_b),\\
 (1-p)\chi'
 &=\left[\sum_{b\sim o}
   \E_p(s_o s_b;\ C_o\ne C_b)\right].
 \end{aligned}
\end{equation}
\end{lemma}
\begin{proof}
In a finite induced domain, differentiate the connection event by
closed pivotals. A closed pivotal edge has a unique orientation
$(a,b)$ for which $x\in C_a$ and $y\in C_b$; these two clusters are
distinct. This proves the finite-domain first identity with the
factor $1-p$ and no extra factor of two.

Fix a compact parameter interval contained in $(0,p_1]$ with
$p_1<\pc$. The summand is bounded, by disjoint witnesses, by
$T_{p_1}(x,a)T_{p_1}(b,y)$. Its sum over $a,b,y$ is at most
$d\chi_{\max}(p_1)^2$. Under exhaustion, connection events converge,
and distinctness indicators converge as well: if two vertices are
connected in the full graph, a finite path eventually witnesses this.
Dominated convergence therefore passes the integrated Russo identity
to the full graph.

For each representative $o_i$, its row is thus the integral of an
$\ell^1(V)$-valued integrand dominated coordinatewise by the fixed
summable row of $(1-p_1)^{-1}T_{p_1}\Adj T_{p_1}$. Lebesgue
differentiation gives row differentiation in $\ell^1$ almost
everywhere. One can see this directly by differentiating a finite
set of coordinates and bounding the remaining tails by the same
summable majorant. There are finitely many row types. Invariance,
symmetry, and the row/column sum bound consequently upgrade these
statements to local absolute continuity and differentiation in
operator norm. Summing the first identity over $y$ and then applying
\eqref{eq:s02} to move the root from $x$ to $a$ gives the second
identity. All sums are justified by the same domination.
\end{proof}

The next connection-operator comparison is
\cite[Lemma~2.4]{Hutchcroft2019}. We include the argument because it is
used both before and after the critical bound is proved.

\begin{lemma}[Sprinkling]\label{pre:sprinkling}
For $0\le p,a\le1$,
\begin{equation}\label{eq:s05}
 T_{p+(1-p)a}
 \le\sum_{j\ge0}a^jT_p(\Adj T_p)^j
 \qquad\text{entrywise}.
\end{equation}
If $T_p$ is bounded and $ad\norm{T_p}<1$, the kernel on the left is
bounded. In particular, boundedness of $T_\pc$ implies
$\pc<\ptwo$.
\end{lemma}
\begin{proof}
Take independent fields of parameters $p$ and $a$; their union has
parameter $p+(1-p)a$. A finite simple open path in the union can be
split at the distinct ordered edges for which it uses the second
field. The intervening first-field legs have disjoint open witnesses.
For prescribed legs and $j$ prescribed second-field edges, independence
and the disjoint-witness inequality give the corresponding product
of $j+1$ connection probabilities times $a^j$. Summing over all
intermediate vertices and $j$ yields \eqref{eq:s05}; extra paths and
repetitions in that sum only increase the bound. Since
$\norm{\Adj}\le d$, the right-hand series converges in operator norm
when $ad\norm{T_p}<1$. Entrywise domination gives the assertion.
Finally $\pc<1$ permits a positive $a$ satisfying this condition at
$p=\pc$.
\end{proof}

\section{Critical clusters remain controlled after deletions}
\label{good:section}

The local maps constructed in the next section attach pieces cut from
critical clusters. We need a property of the original cluster that
controls the separating bridges of every such piece, even after an
arbitrary vertex deletion. This section proves that a capped critical
cluster fails this property only with stretched-exponentially small
probability. All percolation fields in this section have parameter
$\pc$.

\begin{definition}[Pivotal distance and good clusters]
\label{good:definition}
A bridge of a connected open graph is an open edge whose deletion
disconnects that graph. The \emph{pivotal distance} between two vertices
is the number of bridges separating them. Equivalently, delete all
bridges, contract the resulting connected pieces, and measure distance
in the resulting tree. Write $\operatorname{diam}_{\mathrm{piv}}$ for
the maximum pivotal distance and
$\operatorname{rad}_{\mathrm{piv}}(H,x)$ for its maximum from $x$.

Fix $0<\epsilon<1/2$ and set $s=1/2+\epsilon$. A finite open cluster
$C$ with $|C|\le M$ is \emph{good for cap $M$} if the following holds
for every set $A\subseteq C$. For every open component $B$ after
deleting $A$, let $K$ be the number of original open edges from $B$
to $A$. Whenever $K\ge1$,
\[
 \operatorname{diam}_{\mathrm{piv}}(B)\le K M^s.
\]
The pivotal distance in this condition is computed inside $B$.
\end{definition}

\begin{proposition}[Good critical clusters]\label{good:good-clusters}
For each $\epsilon\in(0,1/2)$, uniformly in $x$ and for all sufficiently
large $M$,
\begin{equation}\label{eq:s06}
 \Prob_\pc(s_x\le M,\ C_x\text{ is not good for cap }M)
 \le e^{-M^\epsilon}.
\end{equation}
The event that a cluster has size at most $M$ and is good for cap $M$
is determined by the percolation bits in a deterministic finite-radius
ball about its root.
\end{proposition}

We prove the proposition through a marked-cluster estimate. Independent
vertex marks, often called a ghost field, are a standard device in
percolation differential inequalities \cite{AizenmanBarsky1987}.
The exploration-score method is closely related to that used for the
two-ghost inequality in \cite[Section~3]{Hutchcroft2020Locality};
the deletion-uniform estimate below requires the additional inverse
comparison at the end of this section. In a
deterministic induced domain $D$ containing $x$, mark vertices
independently with probability $h\in(0,1/2]$, independently of the
edges. Let $E=E(D,x,h)$ be the event that $C_x^D$ contains a mark.
An edge is pivotal for $E$ if changing its state changes $\one_E$,
with the marks held fixed. Let $D_p$ denote the number of open
pivotal edges; this notation is confined to the present section and
the percolation parameter is $p=\pc$.

\begin{lemma}[Marked pivotal moments]\label{good:marked-pivotals}
There is a constant $C$, depending only on $d$ and $\pc$, such that,
with
$A_h=C\sqrt{\log(1/h)/h}\ge1$, for every induced domain $D$, every
$x\in D$, every integer $i\ge0$, and every $r\ge0$,
\begin{equation}\label{eq:s07}
 \begin{aligned}
 \E\left[\binom{D_p}{i}\one_E\right]
 &\le A_h^i\Prob(E),\\
 \Prob(E,\ D_p\ge r)
 &\le2\Prob(E)e^{-r/(4A_h)}.
 \end{aligned}
\end{equation}
\end{lemma}
\begin{proof}
\emph{The first moment.}
Explore the cluster in a fixed mark-independent order, testing the
mark of each newly discovered vertex immediately and stopping at the
first mark or when the cluster is exhausted. No edge is queried twice.
Critical finiteness from \cref{pre:critical-finiteness} makes this a
terminating decision procedure, even in an infinite domain.
Let $T$ be its number of edge queries. Conditional on the complete
percolation configuration, the discovery order is fixed. The first
mark index is geometric with parameter $h$, stopped at $|C_x^D|$;
conditional on success, it is that geometric variable conditioned to
be at most $|C_x^D|$. In both cases its mean is at most $1/h$.
There are at most $d$ edge queries per discovered vertex, so
\[
 \E T\le d/h,\qquad \E[T\mid E]\le d/h.
\]

Let $Z=\sum_{e\text{ queried}}(\omega_e-\pc)$. The indicator $Q_e$
that edge $e$ is queried is measurable without its own bit, and every
pivotal edge must be queried by a procedure deciding $E$. Conditioning
on all other edge bits and marks, and then averaging, gives
\[
 \E[(\omega_e-\pc)Q_e\one_E]
 =\pc(1-\pc)\Prob(e\text{ is pivotal}).
\]
The identity remains valid after summation, since $\E T<\infty$.
An open pivotal contributes to $D_p\one_E$, and its bit is independent
of its pivotality. Therefore
\begin{equation}\label{good:score-identity}
 \E[Z\one_E]=(1-\pc)\E[D_p\one_E].
\end{equation}

Each queried bit is fresh conditional on the previous transcript.
For a centered Bernoulli variable, the logarithmic moment-generating
function is at most $\lambda^2/8$: its second derivative is a
Bernoulli variance, at most $1/4$. Successive conditional
expectations at the bounded stopping times $T\wedge n$, followed by
Fatou's lemma, give
\[
 \E\exp(\lambda Z-\lambda^2T/8)\le1.
\]
Put $H=\Prob(E)\ge h$, since a mark at $x$ suffices. Conditional
Jensen's inequality yields
\[
 \lambda\E[Z\mid E]
 \le\log(1/H)+\frac{\lambda^2}{8}\E[T\mid E]
 \le\log(1/h)+\frac{\lambda^2d}{8h}.
\]
Taking $\lambda=\sqrt{8h\log(1/h)/d}$ and using
\eqref{good:score-identity}, we obtain
\[
 \E[D_p\mid E]
 \le\frac{\sqrt{d/2}}{1-\pc}
       \sqrt{\frac{\log(1/h)}h}.
\]
This proves the first-moment assertion, uniformly over domains.

\emph{Higher binomial moments.}
On $E$, each open pivotal is a bridge separating $x$ from all marks
in its cluster. These pivotals are linearly ordered from $x$ toward
the marks. Define
\[
 b_i(D,x)=\E\left[\binom{D_p}{i}\one_E\right],
 \qquad b_0(D,x)=\Prob(E).
\]
For an oriented edge $\vec e=(a,b)$, let $C_{x,\vec e}^D$ be the
root cluster when $e$ is forced closed. Decompose a selected set of
$i$ pivotals according to its first selected edge, oriented away
from $x$. The resulting exact identity is
\begin{equation}\label{good:pivotal-recursion}
 b_i(D,x)=\pc\sum_{\vec e=(a,b)\in E(D)}
 \sum_{\substack{W\subset D\text{ finite}\\x,a\in W, b\notin W}}
 \Prob_\pc(C_{x,\vec e}^D=W)(1-h)^{|W|}
 b_{i-1}(D\setminus W,b).
\end{equation}
To verify the conditional law in this formula, force $e$ closed and
expose the exact finite root cluster $W$. This reveals internal and
boundary information of $W$ but no edges internal to $D\setminus W$.
The event $C_{x,\vec e}^D=W$ ignores the actual bit of $e$.
Require that $W$ contain no marks, accounting for $(1-h)^{|W|}$.
Opening $e$, with independent probability $\pc$, makes it the sole
open contact from $W$ to its exterior. It is then pivotal exactly
when $b$ connects to a mark in $D\setminus W$. All pivotals after
it are precisely those of that exterior event. The exterior edges
and marks retain their independent laws. Every selected subset is
counted once: if the total number of ordered pivotals is $n$, the
count is $\sum_{j=1}^n\binom{n-j}{i-1}=\binom ni$.
All sums are nonnegative, so Tonelli's theorem justifies the
decomposition without a priori moment assumptions.

Uniform induction in $D,x$ now gives
\[
 b_i(D,x)\le A_h^{i-1}b_1(D,x)\le A_h^ib_0(D,x).
\]
Finally expand $(1+(2A_h)^{-1})^{D_p}$ and sum these binomial
moments to get
\[
 \E\left[\left(1+\frac1{2A_h}\right)^{D_p}\one_E\right]
 \le2\Prob(E).
\]
Since $\log(1+(2A_h)^{-1})\ge(4A_h)^{-1}$ for $A_h\ge1$,
Markov's inequality proves the second part of \eqref{eq:s07}.
\end{proof}

\begin{lemma}[A capped pivotal-radius tail]\label{good:radius-tail}
For the fixed $\epsilon\in(0,1/2)$, there is $c>0$ such that, uniformly
in $x$, all sufficiently large $M$, and integers
$1\le K\le M^{1-s}$,
\begin{equation}\label{eq:s08}
 \Prob_\pc\left(s_x\le M,
     \operatorname{rad}_{\mathrm{piv}}(C_x,x)\ge KM^s/2\right)
 \le\exp\left(-\frac{cK^2M^{2\epsilon}}{\log M}\right).
\end{equation}
\end{lemma}
\begin{proof}
On the event on the left, choose a witness vertex $y$ by a fixed
ordering. Independently mark vertices with probability $h$. The
conditional probability that $y$ is the sole mark in $C_x$ is at
least $h(1-h)^M$. In that case every bridge separating $x$ from $y$
is pivotal for the marked event. By \eqref{eq:s07}, with $R=KM^s/2$,
\[
 \Prob_\pc(s_x\le M,\operatorname{rad}_{\mathrm{piv}}\ge R)
 \le\frac{2}{h(1-h)^M}\exp\left(-\frac{R}{4A_h}\right).
\]
Set
\[
 X=\frac{K^2M^{2\epsilon}}{\log M},\qquad
 h=\frac{X}{bM}
   =\frac{K^2M^{-1+2\epsilon}}{b\log M},
\]
where $b$ will be a sufficiently large fixed constant. The allowed
range of $K$ gives $h\le1/(b\log M)\le1/2$. For all sufficiently
large $M$, uniformly in that range,
$\log(1/h)\le2\log M$ and
\[
 \frac{R}{4A_h}\ge\frac{X}{8C\sqrt{2b}}.
\]
Using $-\log(1-h)\le2h$, the logarithm of the preceding probability
bound is therefore at most
\[
 O(\log M)+\frac{2X}{b}-\frac{X}{8C\sqrt{2b}}.
\]
First choose $b$ large enough that the last two terms have a negative
sum bounded above by $-c_1X$. Then increase $M$: the minimum possible
$X$ is $M^{2\epsilon}/\log M$, which dominates $\log M$. This proves
\eqref{eq:s08} with a fixed positive $c$.
\end{proof}

\begin{proof}[Proof of \cref{good:good-clusters}]
Given a capped cluster violating goodness, choose a deleted set $A$
and a remaining component $B$ witnessing the violation, by a fixed
finite ordering. Let $K\ge1$ be the number of original open edges
from $B$ to $A$. Since
$\operatorname{diam}_{\mathrm{piv}}(B)\le|B|-1<M$, a violation has
$K\le M^{1-s}$.

Close its $K$ open contacts, except for one retained contact if
$x\notin B$: in that case take a simple open path from $x$ to $B$
and retain its first edge entering $B$. The path before that contact
does not use any edge being closed, so the modified root cluster
$C'_x$ still contains $B$. If $x\in B$, then $C'_x=B$. In either
case $B$ has at most one open attachment to the rest of $C'_x$.
An internal bridge of $B$ cannot be bypassed outside $B$, because
such a bypass would need two attachments. Thus its bridge tree
embeds into the bridge tree of $C'_x$ with its distances preserved.
The triangle inequality shows that at least one endpoint of a
diameter of $B$ is at pivotal distance greater than $KM^s/2$ from
$x$. Also $|C'_x|\le|C_x|\le M$. Hence the modified configuration
belongs to the event in \eqref{eq:s08}.

It remains to bound the cost of this configuration change; no
independence of its input-dependent choices is needed. At a fixed
output, every changed edge is incident to $B\subseteq C'_x$, so
there are at most $dM$ candidate edges. For fixed $K$, specifying
the subset of at most $K$ changed edges reconstructs the input by
reopening them. The number of possible inputs is at most
\[
 \sum_{j=0}^K\binom{\lfloor dM\rfloor}{j}\le(K+1)(dM)^K,
\]
and their probability ratio to the output is at most
$\max(1,\pc/(1-\pc))^K$. This is a comparison on a finite product:
on the cap event every cluster vertex lies within distance $M-1$
of $x$, and every incident edge lies in a deterministic slightly
larger ball. Condition on all bits outside that ball and apply the
pointwise finite-product comparison inside it.

It follows from \eqref{eq:s08} that the contribution for fixed $K$
is at most
\[
 \exp\left(C_1K\log M-
                 \frac{cK^2M^{2\epsilon}}{\log M}\right).
\]
Uniformly for $K\ge1$, the negative term absorbs the positive one,
because $M^{2\epsilon}/(\log M)^2\to\infty$. Summing over the at
most $M$ possible values of $K$ gives, for large $M$,
\[
 M\exp\left(-\frac{cM^{2\epsilon}}{2\log M}\right)
 \le e^{-M^\epsilon},
\]
as claimed. Finally, explore the root cluster until it is exhausted
or more than $M$ vertices are found. This uses only a deterministic
ball of radius at most $\lceil M\rceil+1$. In the first case its
internal edges and closed boundary are known, and one can test
every subset deletion of the finite cluster. In the second case
the capped event fails. This proves locality.
\end{proof}

\section{Amplifying local endpoint maps}\label{amp:section}

Throughout this section percolation has parameter $\pc$, and the automorphism
group is unimodular.  We use the weighted transport identity
\eqref{eq:s02}, the walk bound \eqref{eq:s03}, critical finiteness, and the
goodness estimate \eqref{eq:s06}.  The purpose of the construction is to turn
a heavy cluster-size tail into a local rule whose endpoint spreads very
widely relative to its probability of success and its cluster-size cap.
The next section will show that such rules contradict the heavy tail.

Suppose that
\begin{equation}\label{eq:s09}
 [\E_{\pc}s_o^\alpha]=\infty
 \qquad\text{for some }0<\alpha<\tfrac12.
\end{equation}
Choose $\beta>\alpha$ and $\epsilon>0$ so that, with
$s=\tfrac12+\epsilon$, we have $\beta+s<1$.  Goodness below always uses this
choice of $s$.

\subsection{One orbit and a sequence of scales}\label{amp:scales-section}

\begin{lemma}[Selection of an orbit]\label{amp:scales}
There are a vertex orbit $O$, a constant $c>0$, and an unbounded sequence
$\mathcal N$ of dyadic integers such that, for $n\in\mathcal N$ and
$\pi=[\Prob_{\pc}(n\le s_o<2n)]$,
\begin{equation}\label{eq:s10}
 \begin{split}
 \pi&\ge n^{-\beta},\\
 [\Prob_{\pc}(n/2\le s_o<2n)]&\le(1+2^\beta)\pi,\\
 \E_{\pc}\bigl[|C_x\cap O|\one_{\{n\le s_x<2n\}}\bigr]
     &\ge c n\pi\qquad(x\in O).
 \end{split}
\end{equation}
The orbit and the sequence can be kept fixed throughout the amplification.
\end{lemma}

\begin{proof}
Put $\pi_j=[\Prob_{\pc}(2^j\le s_o<2^{j+1})]$.  The sequence
$2^{\beta j}\pi_j$ is unbounded: a uniform bound would make
$\sum_j2^{\alpha j}\pi_j$ finite, contrary to \eqref{eq:s09}.
Choose its successive record indices, omitting an initial finite segment
so that their record values are at least one.  At a record index $j$,
$\pi_{j-1}\le2^\beta\pi_j$, proving the first two assertions.

Let $h$ be the number of vertex orbits.  Call a vertex of a finite cluster
qualifying if its orbit occupies at least $1/h$ of the cluster.  At least
$1/h$ of the vertices of every cluster qualify.  Transport, uniformly
within each finite cluster in the chosen bin, gives
\[
 [\Prob_{\pc}(n\le s_o<2n,\ o\text{ qualifies})]
 =\left[\E_{\pc}\left[
   \one_{\{n\le s_o<2n\}}
   \frac{\#\{v\in C_o:v\text{ qualifies}\}}{s_o}\right]\right]
 \ge\frac\pi h.
\]
For at least one type $i$, its contribution to the left side is at least
$\pi/h^2$.  At a representative $x$ of this type,
\[
 \E_{\pc}\bigl[|C_x\cap O_i|\one_{\{n\le s_x<2n\}}\bigr]
 \ge \frac n h\,
      \Prob_{\pc}(n\le s_x<2n,\ x\text{ qualifies})
 \ge \frac{n\pi}{h^3\sigma_i}.
\]
One of the finitely many types is selected along an unbounded subsequence
of record scales.  Keep that subsequence and write $O$ for its orbit.
\end{proof}

\begin{definition}[Local endpoint map]\label{amp:map}
A local endpoint map on $O$ is a randomized rule at each $x\in O$ that
reads percolation bits in a deterministic finite ball about $x$ and uses
auxiliary randomness independent of those bits.  Its conditional rules
are covariant under automorphisms.  The rule either fails or outputs
$y\in C_x\cap O$.  A map with integer cap $M$ is required, on every successful
input pattern, to have $s_x\le M$ and $C_x$ good for cap $M$.

Write $u>0$ for its success probability and
\[
 \begin{gathered}
 K(x,y)=\Prob(\text{output }y\mid\text{success at }x),
 \qquad r=\norm K,\\
 U=\log(1/u),\quad m=\log M,\quad H=\log(1/r).
 \end{gathered}
\]
All probabilities here include the auxiliary randomness, and the operator
norm is on counting-measure $\ell^2(O)$.
\end{definition}

Invariance makes $u$ independent of $x\in O$.  The kernel $K$ has row sums
one; transport restricted to $O\times O$ gives column sums one as well.
Thus $0<r\le1$ by the row and column sum bound.  Symmetry of $K$ is neither
assumed nor needed.

\begin{proposition}[Efficient endpoint maps]\label{amp:efficient-maps}
Under \eqref{eq:s09}, for every $\eta>0$ there is a local endpoint map on
the orbit $O$ of \cref{amp:scales}, with $M\ge2$ and $H>1$, such that
\[
 \frac{U+m}{H}<\eta.
\]
Moreover the maps can be chosen so that $H$ tends to infinity.  Their
radii are finite, but need not be bounded uniformly over the construction.
\end{proposition}

We prove the proposition by constructing a starting map, establishing
one amplification step, and then specifying the order of all parameters.

\begin{lemma}[Starting map]\label{amp:initial}
There is a local endpoint map on $O$ with $H>1$, with cap as large as
prescribed.
\end{lemma}

\begin{proof}
From $x\in O$ sample a length-$\ell$ walk with kernel $P$, independently of
percolation.  Require its endpoint to lie in $O$ and all distinct edges
traversed by the walk to be open.  If $e$ is the number of distinct
traversed edges, accept also an independent coin of probability
$\pc^{\ell-e}$.  Holding steps do not count as edges.  Before imposing a
cluster cap, the resulting endpoint mass kernel is exactly
$\pc^\ell P^\ell|_{O\times O}$.  The finite type chain gives, for all large
$\ell$, a row sum at least $c_1\pc^\ell$, with $c_1>0$ independent of
$\ell$.  Its norm is at most $\pc^\ell\rho^\ell$.

Fix $\ell$ large enough that $2\rho^\ell/c_1<e^{-1}$.  Critical finiteness
and \eqref{eq:s06} allow us to choose any sufficiently large $M$ and reject
unless $s_x\le M$ and $C_x$ is good for cap $M$, losing at most half the
preceding lower bound on the row sum.  Entrywise domination by the uncut
mass kernel gives $r\le2\rho^\ell/c_1<e^{-1}$.  The capped cluster and its
goodness can be tested in a deterministic finite ball, so this is a local
map.
\end{proof}

\subsection{Splicing independent clusters}\label{amp:splicing-section}

Fix an old map with parameters $(u,M,r)$, an integer $k\ge1$, and a scale
$n\in\mathcal N$ such that $kM\le n/2$.  Constants in what follows may
depend on the graph and the fixed choices $\beta,\epsilon,O$, but not on
$n,k$ or the old map unless stated otherwise.

All finite choices can be made covariantly using independent continuous
priorities on vertices or edges.  More explicitly, keep three independent
families of seeds: the old-map seeds, priorities ordering new contact
lists, and a uniform random variable $V_i\in[0,1]$ for the contact choice
at each possible index $1\le i\le k$.  Ties in priorities have probability
zero and may be rejected.  An automorphism permutes the seed indices and
preserves their joint law, so this realization does not require a free
action or a choice of an automorphism carrying one root to another.
A rule with already integrated endpoint probabilities is likewise
realized by a categorical draw using its old-map seed.  Once those seeds
are fixed, old rules can be replayed deterministically.

For fixed $x,z\in O$ perform the following experiment.  In a distinguished
independent field $D$, require
$A=C_x(D)$ to satisfy $n\le|A|<2n$.  Independently, start at $z$ and take
$k$ old-map steps, using a new independent percolation field and new seeds
for each step.  Give the experiment mass $1/n$ if every step succeeds and
the last endpoint lies in $A$, and zero otherwise.  A successful chain,
listed in reverse order, has vertices
\[
 y_0\in A,\quad y_1,\ldots,y_k=z.
\]
Let $E_i$ be the field used between $y_i$ and $y_{i-1}$, and let $C_i$ be
its whole open cluster containing these vertices.  Each $C_i$ is good for
cap $M$ and has at most $M$ vertices.  This reversal is only a way of
listing a realized chain; it makes no reversibility assertion about $K$.

Starting with $S=A$, apply the following switching procedure until
$z\in S$.  Select the largest still eligible index $i$ for which
$C_i\cap S\ne\varnothing$; after this selection only indices larger than
$i$ remain eligible.  Let $B_i$ be the component of $y_i$ in the open
subgraph of $E_i$ induced by $V\setminus S$.  Among its $K_i$ open
$E_i$-edges to $S$, choose one edge $e_i$ uniformly.  Exchange the $D$ and
$E_i$ bits on the mask
\[
 \{e_i\}\ \cup {}
 \{e:e\text{ is incident to }B_i
             \text{ and both endpoints of }e\text{ lie outside }S\}.
\]
The field $D$ in this instruction is its current value.

\begin{lemma}[Geometry of the switching procedure]\label{amp:splicing}
Every step is well defined.  Its effect on the distinguished root cluster
is exactly $S\mapsto S\cup B_i$.  The selected indices increase, and the
inserted pieces form a chain, each attached to the preceding piece by
one open bridge $e_i$; the first piece is attached to $A$.
Every selected bridge separates $x$ from $z$ in the final distinguished
field $D'$, and
\begin{equation}\label{eq:s11}
 z\in C_x(D'),\qquad s_x(D')\le M'=2n+kM.
\end{equation}
\end{lemma}

\begin{proof}
Initially an eligible intersecting cluster exists because $y_0\in A$.
After a selection at $i$, the new set contains $y_i$, so if $i<k$ and
$z$ has not been reached, $C_{i+1}$ supplies another eligible intersection.
The selected $y_i$ lies outside the current $S$: if $i<k$ and $y_i\in S$,
then $C_{i+1}$ would intersect $S$, contradicting maximality; if $i=k$,
this would contradict the stopping condition.  Since $C_i$ is connected
and meets $S$, the component $B_i$ has $K_i\ge1$ open contacts with $S$.

Before the exchange, all $D$-edges leaving $S$ are closed.  The exchange
imports the internal open connections of $B_i$ and its closed boundary
to $V\setminus(S\cup B_i)$ from $E_i$, and opens the chosen contact
$e_i$.  Every other contact to $S$ keeps its closed $D$ bit.  No edge
internal to $S$ changes.  Hence the new root cluster is exactly
$S\cup B_i$, with $e_i$ its sole open contact between the two parts.

Every index larger than the selected $i$ missed the entire pre-switch
$S$.  Consequently the next selected cluster can contact the enlarged
root cluster only through $B_i$.  Applying this observation at each step
shows that later pieces attach to their immediate predecessors and do
not acquire open contacts with earlier pieces.  Thus all the seams
remain bridges, and, when the procedure stops, separate $x$ from
$z$ in their natural order.  There are at most $k$ inserted pieces,
each of size at most $M$, proving \eqref{eq:s11}.
\end{proof}

\Cref{fig:splicing} illustrates the chain of inserted pieces.  The
separating seams will let us recover the original fields from the
switched fields once the seam locations have been recorded.
\begin{figure}[tbp]
\centering
\begin{tikzpicture}[x=1cm,y=1cm,font=\small,
  piece/.style={draw,rounded corners=8pt,fill=blue!5,minimum width=1.65cm,minimum height=1.05cm},
  seam/.style={very thick,blue!65!black},dot/.style={circle,fill,inner sep=1.4pt}]
\node[piece,minimum width=2cm] (a) at (0.3,0) {$A$};
\node[piece] (b1) at (3.1,0) {$B_{i_1}$};
\node[piece] (b2) at (5.8,0) {$B_{i_2}$};
\node[piece] (bt) at (8.5,0) {$B_{i_t}$};
\draw[seam] (a.east)--node[above] {$e_{i_1}$} (b1.west);
\draw[seam] (b1.east)--node[above] {$e_{i_2}$} (b2.west);
\draw[seam,densely dotted] (b2.east)--node[above] {$\cdots$} (bt.west);
\node[dot,label=below:$x$] at (-0.2,0.10) {};
\node[dot,label=below:$z$] at (8.9,0.10) {};
\node[align=center,text width=10.7cm] at (4.3,-1.25)
{Each selected piece is attached by one open seam.\\
The seams occur in order among the bridges separating $x$ and $z$.};
\end{tikzpicture}
\caption{The chain produced by the switching procedure. Cutting the
selected seams recovers the inserted pieces and the initial cluster
$A$. The drawing shows vertex sets and selected open seams, with
internal edges and other incident edges omitted.}
\label{fig:splicing}
\end{figure}


\subsection{Coding an inverse and defining its load}\label{amp:coding-section}

\begin{lemma}[Inverse codes and contact probabilities]\label{amp:coding}
Fix the output fields, $x,z$, the old-map seeds, and the priorities
ordering contact lists.  A valid switching input can be specified by the
selected indices and positive gaps locating the selected seams in the
ordered list of bridges separating $z$ from $x$, read from $z$ toward
$x$.  Each fixed code determines at most one input bit tuple, without
using the fresh contact uniforms.  After integrating those uniforms, the
total mass of valid codes is at most
\begin{equation}\label{eq:s12}
 W^k,\qquad
 W=1+\sum_{l=1}^{M}\min\{1,2M^s/l\}
       \le C M^s\log M.
\end{equation}
\end{lemma}

\begin{proof}
List the selected indices in decreasing order, hence their seams in the
order encountered from $z$ toward $x$.  The code records the number of
separating bridges up to each next seam, counting that seam.  Cutting
the indicated seams in the final root cluster recovers the inserted
pieces and $A$, since \cref{amp:splicing} gives exactly a chain of singly
attached connected pieces.  Thus it determines every pre-switch set
$S$, every $B_i$, and every exchange mask.  Undo these masks in reverse
chronological order.  The candidate input tuple is now fixed, and the
old rules can be replayed from $z$ to check all their successes, their
endpoints, and the prescribed selection rule.  An invalid code is simply
discarded.

Let $l_i$ be the gap ending at seam $e_i$.  Other than that seam, its
bridges lie inside $B_i$, between the attachment point toward $z$ and
the attachment point toward $x$.  Therefore $l_i\le|B_i|\le M$.
Deleting the pre-switch set $S$ from the original good cluster $C_i$
leaves $B_i$ with $K_i$ original open contacts.  Goodness bounds its
pivotal diameter by $K_iM^s$.  Since no later piece supplies a bypass,
\[
 l_i\le1+K_iM^s,
 \qquad
 \frac1{K_i}\le\min\{1,2M^s/l_i\}.
\]
For the second inequality use $K_i\ge1$ if $l_i=1$, and
$l_i-1\ge l_i/2$ otherwise.

The reconstruction just given has not consulted any $V_i$.  In the
reconstructed input, each contact list and its priority order are fixed.
If its prescribed seam has rank $a_i$ among $K_i$ contacts, the fresh
uniform must lie in
$[(a_i-1)/K_i,a_i/K_i)$.  These intervals are determined without any of
the fresh uniforms, so their product Lebesgue measure is
$\prod_iK_i^{-1}$.  Any further validity restrictions only decrease it.
The priorities are integrated once under their probability law, rather
than counted as additional inverse choices.  Summing over a skipped
index or a gap $1\le l\le M$ at each index proves the bound $W^k$.
Finally the harmonic sum gives $W\le C M^s\log M$ for $M\ge2$.
\end{proof}

For clarity, all measure comparisons here can be made on finite bit
spaces.  In a successful experiment the chain connects $z$ by at most
$kM$ edges to a vertex of the connected set $A$, so
$\dist(x,z)\le2n+kM$.  A capped whole-cluster test exposes the boundary
or stops as soon as the cap is exceeded.  The old rule has a fixed query
radius, and its successful steps move at most $M$.  Consequently a
deterministic larger ball about $x$ contains every query and every
exchange mask for every successful experiment.  There are finitely many
candidate endpoints and codes.  When checking an inverse, stop at the
first failed capped test or failed old step; this uses a deterministic
larger ball as well.  Conditioning on all bits outside this ball reduces
every comparison to finitely many Bernoulli coordinates.

At fixed seeds, the exchanges preserve the total number of open bits
among equally distributed fields.  Thus each input bit tuple and its
output have exactly the same product probability, even though the
exchange masks depend on the input.  The injective inverse with a fixed
code therefore gives the following exact density description.  For a
single reference iid distinguished field $D'$, let
$L_x(z,D')$ be $1/n$ times the sum of valid inverse-code indicators,
integrated over reference iid auxiliary output fields and all independent
seeds.  The preceding finite comparison says that this is the density
of the switched mass at $z$ relative to the law of $D'$.  In particular,
\[
 0\le L_x(z,D')\le W^k/n.
\]
After the auxiliary integrations these are covariant local functions
of $D'$, supported on \eqref{eq:s11}.  Put
$L=\sum_{z\in O}L_x(z,D')$ and $b_0=\E L$.
The reference product law in this definition is not a claim that the
actual output of a switch has independent fields, or that its contact
coins remain independent after conditioning on that output.

\begin{lemma}[Mean load]\label{amp:mean}
With $F(x,z)=\E L_x(z,D')$, one has
\begin{equation}\label{eq:s13}
 \begin{split}
 F(x,z)&=\frac{u^k}{n}\sum_{y\in O}
  \Prob_{\pc}(y\in C_x,\ n\le s_x<2n)K^k(z,y),\\
 c\pi u^k&\le b_0\le C\pi u^k,
 \qquad \norm{F/b_0}\le r^k.
 \end{split}
\end{equation}
\end{lemma}

\begin{proof}
The density comparison preserves total mass, so the mean is computed
in the original experiment.  The $k$ independent old steps have
unconditional endpoint kernel $u^kK^k$, giving the displayed formula.
Let
\[
 B_n(x,y)=\Prob_{\pc}(y\in C_x,\ n\le s_x<2n),\qquad x,y\in O.
\]
This is symmetric: on $y\in C_x$, the two root clusters coincide.
Its row sum $b_n$ is constant on $O$, and by \eqref{eq:s10} and
$\Prob_{\pc}(n\le s_x<2n)\le\pi/\sigma_O$,
\[
 cn\pi\le b_n\le2n\pi/\sigma_O.
\]
Here $\sigma_O$ is the type weight of $O$.  Since $K^k$ has column sums
one, $b_0=u^kb_n/n$.  Moreover
$F/b_0=(B_n/b_n)(K^k)^*$, where $*$ denotes transpose.  The symmetric
stochastic kernel $B_n/b_n$ has norm at most one, proving the norm bound.
\end{proof}

\subsection{The second moment}\label{amp:second-moment-section}

The mean kernel already has norm at most $r^k$ after normalization.
To turn $L_x(z,D')$ into output probabilities, however, we must control
the fields on which its total load $L$ is large.  The next estimate will
let us truncate those fields while retaining a fixed fraction of the
mean.  The distinction between an underlying graph edge and an open
edge is essential in its proof.

\begin{lemma}[Second moment of the load]\label{amp:second-moment}
There is a fixed polynomial $Q(k,M)\ge1$ such that
\begin{equation}\label{eq:s14}
 \E L^2\le Q(k,M)\pi u^{2k}
 \qquad\text{whenever } kM\le n/2
               \text{ and } n\ge(W/u)^k.
\end{equation}
\end{lemma}

\begin{proof}
Write $v=kM$.  Expand $L^2$ using a common reference distinguished field
$D'$ and two independent families of auxiliary output fields and seeds.
Thus an individual summand consists of two valid coded inverses with
endpoints $z_1,z_2$.  We divide the summands into three classes.

\paragraph{At least one inverse makes no switch.}
Let $L_0$ be the integrated part of $L$ with no switch.  In this case
$A=C_x(D')$ is in the bin and $z\in A$.  Ignoring any further conditions
and summing the old-chain mass gives
\[
 L_0\le\frac{u^k}{n}\one_{\{n\le|A|<2n\}}
          \sum_{z\in A\cap O}\sum_{y\in A\cap O}K^k(z,y)
 \le 2u^k,
\]
by the column sums.  Summands with at least one such factor contribute
at most $2\E[L_0L]\le4u^kb_0\le C\pi u^{2k}$.

\paragraph{The two hanging components are nested or adjacent.}
For a nonempty switching inverse, let $B$ be the component beyond its
first seam, viewed from $x$ in $C_x(D')$.  By \cref{amp:splicing}, $B$ is
exactly the union of the inserted pieces, contains its endpoint $z$,
and has at most $v$ vertices.  Two such hanging components in the same
rooted bridge tree are either nested or disjoint.

Fix the first inverse and its $B^1,z_1$.  There are at most
\[
 N_v=v+v^2+dv^3
\]
possible second endpoints among pairs for which the hanging components
are nested or joined by an underlying graph edge.  Indeed, if
$B^2\subseteq B^1$, there are at most $v$ endpoints.  If
$B^1\subseteq B^2$, the first seam of $B^2$ is one of the last $v$
separating bridges on the route from $x$ to $z_1$.  Each seam determines
one hanging component and supplies at most $v$ endpoints when that
component is small enough.  This gives $v^2$.  Finally, for disjoint
adjacent components, choose a vertex of $B^2$ neighboring $B^1$ in at
most $dv$ ways; its seam is among the last $v$ separating bridges on
the route to that vertex, and each resulting component contains at
most $v$ possible endpoints.  This gives $dv^3$.

For each endpoint, \cref{amp:coding} bounds the integrated second
inverse mass by $W^k/n$.  Integrating the first inverse therefore bounds
this class by
\[
 b_0N_v\frac{W^k}{n}\le C N_v\pi u^{2k},
\]
using $n\ge(W/u)^k$.

\paragraph{The hanging components are disjoint and graph-nonadjacent.}
Now suppose $B^1$ and $B^2$ are disjoint and no edge of the underlying
graph joins them.  Every distinguished-field edge exchanged by inverse
$a$ is incident to $B^a$.  Hence the two collections of masks are
disjoint, including their seams.  Undo both inverses.  In the jointly
restored distinguished field $D_{00}$ the exact root cluster is
\begin{equation}\label{amp:joint-root}
 A_{00}=C_x(D')\setminus(B^1\cup B^2),
 \qquad n-v\le|A_{00}|<2n.
\end{equation}
To verify the first assertion, removing two disjoint hanging components
leaves a connected set containing $x$.  Its contacts with each removed
component are closed by that inverse, and are untouched by the other
inverse.  Its other boundary edges remain closed.  For the size bound,
either single inverse starts from
$A^a=C_x(D')\setminus B^a$ of size in $[n,2n)$, and removing the other
component loses at most $v$ vertices.

The geometry of this joint restoration is shown in
\cref{fig:joint-inverse}.
\begin{figure}[tbp]
\centering
\begin{tikzpicture}[x=1cm,y=1cm,font=\small,
  piece/.style={draw,rounded corners=8pt,fill=blue!5,minimum width=1.65cm,minimum height=1.05cm},
  seam/.style={very thick,blue!65!black},dot/.style={circle,fill,inner sep=1.4pt}]
\node[piece,minimum width=2.2cm] (root) at (0.4,0) {$A_{00}$};
\node[piece,minimum width=2.1cm] (upper) at (4.3,.35) {$B^1$};
\node[piece,minimum width=2.1cm] (lower) at (4.3,-1.1) {$B^2$};
\draw[seam] (root.east) to[bend left=10] node[above] {first seam} (upper.west);
\draw[seam] (root.south east) to[bend right=8] node[below,sloped] {first seam} (lower.west);
\node[dot,label=below:$x$] at (-0.2,-.05) {};
\node[dot,label=right:$z_1$] at (4.75,.35) {};
\node[dot,label=right:$z_2$] at (4.75,-1.1) {};
\draw[decorate,decoration={brace,amplitude=4pt}] (5.7,.87)--(5.7,-1.62);
\node[align=left,anchor=west,text width=3cm] at (6,-.35)
{No underlying graph edge joins $B^1$ to $B^2$. Their exchange masks are disjoint.};
\end{tikzpicture}
\caption{Two hanging components in a common switched root cluster.
Undoing both insertions leaves $A_{00}$. Graph nonadjacency makes
the two collections of exchange masks disjoint. The proof recovers
each insertion from the jointly restored fields. The open seams
shown belong to the switched configuration; internal and other
incident edges are omitted.}
\label{fig:joint-inverse}
\end{figure}


We claim that, at fixed $z_1,z_2$ and all seeds, joint undoing is injective
without retaining either code.  This stronger assertion removes the
factor $W^{2k}$ from this class.  The restored auxiliary fields and
old-map seeds first recover both entire old chains by replay from their
endpoints; the distinguished field is irrelevant to that replay.
Run the first switching procedure virtually from $A_{00}$ with its
restored chain.  Here ignore only the original bin test and the test
that $y_0\in A$; keep the maximal-index selection and stopping rules.
We show inductively that its current set is $S_0=S\setminus B^2$, where
$S$ is the corresponding current set in the original first forward run,
and that the same index, component, contact and mask are selected.

Let $i$ be the original selected index and write its seam as
$e_i=(a_i,b_i)$ with $a_i\in S$ and $b_i\in B_i^1\subseteq B^1$.
Graph nonadjacency forces $a_i\notin B^2$, so the selected old cluster
still meets $S_0$.  Every larger eligible index missed $S$ and therefore
misses $S_0$.  Thus maximality chooses the same $i$.  This remains true
when unused trajectory clusters meet $B^2$, or even when the original
$y_0$ belongs to $B^2$.

Replacing $S$ by $S_0$ admits only the vertices of $B^2$ to the slit
domain.  If the component $B_i^1$ enlarged, an enlarging open path would
have a first vertex in $B^2$.  The preceding part of that path lies in
the old slit domain, hence in $B_i^1$.  Its next edge would join $B^1$
to $B^2$, a contradiction.  Thus the slit component is unchanged.
Its contact list and its exchange mask are also unchanged, since their
only possible differences would be edges between $B_i^1$ and $B^2$.
The fixed priorities and contact uniform choose the same seam.  The
other undoing changed no bit on this mask.  After the exchange the
inductive relation $S_0=S\setminus B^2$ therefore continues to hold.
The stopping tests agree because $z_1\in B^1$ and $z_1\notin B^2$.

The same proof applies to the second virtual replay.  Run the two
replays separately from the jointly restored tuple and combine their
disjoint distinguished-field changes.  This reconstructs the common
output field, both auxiliary output families, the seams, and hence
both codes.  The claimed injection follows.  Also the first selected
seam in each replay meets $A_{00}$; thus each restored successful
trajectory, meaning the union of its old-map whole clusters, hits
$A_{00}$.

The joint undo preserves the probability of the full finite Bernoulli
bit tuple. More explicitly, fix all seeds and let $N$ be the number
of bit coordinates in the finite comparison. A tuple with $a$ open
coordinates has weight $\pc^a(1-\pc)^{N-a}$, and the exchanges preserve
$a$. Preservation of this weight by itself would not control
multiplicity. The replay just proved makes the map from a valid
output bit tuple together with its two codes to the jointly restored
bit tuple injective. Summing its weights therefore gives exactly
the mass of its image, bounded by the mass of the unrestricted
reference experiment. Finally integrate the fixed seeds by Tonelli's
theorem. The injection thus bounds this class by an unrestricted
reference product experiment, weighted by $1/n^2$, with an iid root
field and two independent old-chain families.  Retain only the size
window in \eqref{amp:joint-root}, success of both chains, and the
condition that each trajectory hits the root cluster.  Independence of
the two chain families, conditional on that cluster, is invoked in this
reference experiment only, after the injective comparison.

For a deterministic finite set $A$, let $h_z(A)$ be the probability that
a successful $k$-step old trajectory from $z\in O$ hits $A$.  Then
\begin{equation}\label{amp:trajectory-hit}
 \sum_{z\in O}h_z(A)\le C|A|v u^k.
\end{equation}
Indeed send unit mass from a starting vertex $z\in O$ to every vertex
of its trajectory union on success, and send zero from other types.
The expected outgoing mass at $z$ is at most $v u^k$, because each
fresh step succeeds with probability $u$ and the union has at most $v$
vertices.  Weighted transport bounds the expected incoming mass at
every vertex by $v u^k/\min_i\sigma_i$.  Summing over vertices of $A$
and applying a union bound proves \eqref{amp:trajectory-hit}.

Conditional on the reference root cluster $A$, the summed contribution
of both chains is bounded by the square of \eqref{amp:trajectory-hit}.
Since $v\le n/2$, the contribution of this class is at most
\[
 \frac{C v^2u^{2k}}{n^2}
 \E_{\pc}\bigl[s_x^2; n-v\le s_x<2n\bigr]
 \le C v^2u^{2k}\Prob_{\pc}(n/2\le s_x<2n)
 \le C'v^2\pi u^{2k}.
\]
The final inequality uses \eqref{eq:s10} and the positive type weight
$\sigma_O$.  Combining the three classes proves \eqref{eq:s14}; for
example a sufficiently large graph-dependent multiple of
$1+v+v^2+dv^3$ is an admissible polynomial $Q(k,M)$.
\end{proof}

\subsection{Truncation and iteration}\label{amp:iteration-section}

\begin{lemma}[One amplification step]\label{amp:step}
Fix an old map.  For $n\in\mathcal N$ tending to infinity and
$k=O(\log n)$ with $kM\le n/2$ and $n\ge(W/u)^k$, there is, for all
sufficiently large such $n$, a new local map with parameters satisfying
\begin{equation}\label{eq:s15}
 u'\ge c'\pi/Q(k,M),\qquad r'\le2r^k,
 \qquad M'=2n+kM.
\end{equation}
The threshold for $n$ may depend on the fixed old map.
\end{lemma}

\begin{proof}
Choose a fixed sufficiently large constant $C_0$ and put
$b=C_0Q(k,M)u^k$.  By \eqref{eq:s13} and \eqref{eq:s14}, discarding
fields on which $L>b$ loses mean mass at most
\[
 \E[L;L>b]\le\frac{\E L^2}{b}\le b_0/4.
\]
On the remaining fields, discard also those whose root cluster is not
good for cap $M'$.  The load is already supported on the cap by
\eqref{eq:s11}.  Thus \eqref{eq:s06} bounds this second loss by
$b\exp(-(M')^\epsilon)$.  Relative to $b_0\ge c\pi u^k$, it is at most
\[
 C Q(k,M)n^\beta\exp(-(M')^\epsilon),
\]
which tends to zero when the old map is fixed and $k=O(\log n)$.
For sufficiently large $n$ the second loss is also at most $b_0/4$.

On a retained field output $z$ with probability $L_x(z,D')/b$, and
fail with the remaining probability.  This is legitimate because
$L\le b$.  It is local and covariant, and every successful input has
its endpoint in a good root cluster of size at most $M'$.  If $F'$ is
the retained, unnormalized endpoint mass kernel, its row sum $b_0'$
satisfies $b_0'\ge b_0/2$, and $0\le F'\le F$ entrywise.  Hence
\[
 u'=b_0'/b\ge c'\pi/Q(k,M),\qquad
 r'=\norm{F'/b_0'}\le2\norm{F/b_0}\le2r^k.
\]
The probabilities defining this rule are now fixed numbers.  They can
later be applied to percolation at another parameter without changing
the rule, and its stated physical success conditions still hold.
\end{proof}

\begin{proof}[Proof of \cref{amp:efficient-maps}]
Choose $\xi>0$ so small that
\[
 c_\xi=\frac{\beta+\xi+(s+\xi)(1+\xi)}{1-\xi}<1;
\]
this is possible because $\beta+s<1$.  By \cref{amp:initial}, start with
$H>1$ and $M$ sufficiently large for the estimates below.  For any fixed
old map set
\[
 t=U+(s+\xi)m,
 \qquad a=\left\lceil C M^s\log M/u\right\rceil.
\]
Choose $C$ large enough that $a\ge\max\{2,W/u\}$.  For all sufficiently
large $M$, uniformly in $0<u\le1$, the logarithmic factor and the
rounding cost are absorbed by $\xi\log M$, giving $\log a\le t$.

Only after fixing this old map and $a$ choose $n\in\mathcal N$ large,
and put $k=\lfloor\log n/\log a\rfloor$.  Then $k\to\infty$,
$k=O(\log n)$, and $n\ge a^k\ge(W/u)^k$.  The condition $kM\le n/2$
and all goodness and truncation thresholds eventually hold.  Also
\[
 k\log a\le\log n<(k+1)\log a\le(k+1)t.
\]
By \eqref{eq:s15}, $\pi\ge n^{-\beta}$, and the polynomial dependence
of $Q(k,M)$ on $k$ for fixed $M$, we can make $n$ large enough that
\[
 U'\le k(\beta+\xi)t,\qquad
 m'\le k(1+\xi)t,\qquad
 H'\ge k(1-\xi)H>1.
\]
For example, the first follows from
$U'\le\beta\log n+\log Q(k,M)+O(1)$, where
$\log Q(k,M)=O_M(\log(k+1))$; the second follows from
$M'\le(5/2)n$; the third follows from $r'\le2r^k$.
Consequently, with $t'=U'+(s+\xi)m'$,
\[
 \frac{t'}{H'}\le c_\xi\frac tH.
\]
At each stage choose the new scale still larger if necessary so that
$H'\ge2H$.  This changes none of the preceding inequalities.  Iterate,
keeping the same orbit $O$ and scale sequence $\mathcal N$ throughout.
Then $t/H$ tends to zero and $H$ tends to infinity.  Since
$U+m\le\max\{1,(s+\xi)^{-1}\}\,t$, the desired conclusion follows.
Each individual choice uses a finite radius; no uniform radius is used
in this iteration.
\end{proof}

\section{Critical moments and susceptibility}
\label{mom:section}

We first turn the local maps into a lower bound on the rate of growth of
the susceptibility.  Under the temporary assumption \eqref{eq:s09},
\cref{amp:efficient-maps} supplies maps on a fixed orbit $O$ with
\[
 U=\log(1/u),\qquad m=\log M,\qquad H=\log(1/\norm K),
 \qquad \frac{U+m}{H}\longrightarrow0.
\]
The calculations with susceptibilities and second moments take place
at $p<\pc$, where the preliminary estimates ensure their finiteness.

\begin{lemma}[Maximal susceptibility in an intrinsic ball]
\label{mom:maximal-potential}
For a deterministic induced domain $D$, put
$f_D(v)=\E_p\abs{C_v^D}$ for $v\in D$, and $f_D(v)=0$ otherwise.
Let $\mathcal B^D(v,M)$ be the vertices reachable from $v$ by an open
path in $D$ of length at most the nonnegative integer $M$.  If $v\in D$
and $0<p<\pc$, then
\begin{equation}
 \E_p\max_{z\in\mathcal B^D(v,M)}f_D(z)
 \le \left(1+\frac{1-p}{p}M\right)f_D(v).
 \label{eq:s16}
\end{equation}
\end{lemma}

\begin{proof}
First let $D$ be finite.  Select a maximizer by a fixed tie rule.  For
each $w\in D$, take a percolation configuration $\omega$ and expose the
exact vertex set $B=C_w^D(\omega)$.  Replace every edge with at least one
endpoint in $B$ by an independent Bernoulli($p$) bit, leaving all other
edges unchanged; call the resulting configuration $\widetilde\omega$.
Conditional on $B$, the unchanged exterior edges have product law,
and all replaced bits are fresh.  Consequently $\widetilde\omega$ has
product law independent of $B$.  If $Z$ is the selected maximizer in
$\mathcal B^D(v,M;\widetilde\omega)$, this independence gives
\[
 \sum_{w\in D}\Prob(Z\in B)
 =\sum_{z\in D}\Prob(Z=z)\sum_{w\in D}T_p^D(w,z)
 =\E_p\max_{z\in\mathcal B^D(v,M)}f_D(z).
\]
Here a separate coupling is used for each $w$; only the marginal law of
$Z$ needs to be the same in these couplings.

If $v\notin B$ and $Z\in B$, take the first edge entering $B$ along
a simple $\widetilde\omega$-open path from $v$ to $Z$ of length at most
$M$.  This edge was closed in $\omega$, and its exterior endpoint was
reachable from $v$ in $D\setminus B$ by at most $M-1$ open edges of
$\omega$.  Let $N_w(\omega)$ count boundary edges with this property,
setting it to zero when $v\in B$.  Then
\[
 \Prob(Z\in B)\le T_p^D(v,w)+\E_p N_w.
\]
To bound the last expectation, open one counted, distinguished edge.
In the new configuration it is pivotal for $v\leftrightarrow w$:
it is the only open edge entering the old cluster $B$.  Moreover it is
among the first $M$ bridges separating $v$ from $w$, in their order
from $v$, because a simple path reaches it using at most $M$ edges.
For a fixed output and a distinguished edge, closing that edge
recovers the input.  There are at most $M$ possible distinguished
edges in each connected output.  Comparing the two possible values
of that one bit therefore gives
\[
 \E_p N_w\le \frac{1-p}{p}M\,T_p^D(v,w).
\]
Summation over $w$ proves \eqref{eq:s16}.  For an infinite domain,
exhaust $D$ by finite induced domains containing $v$.  Their
susceptibilities and intrinsic balls increase to the corresponding
objects in $D$.  The maximum also converges monotonically: the
intrinsic $M$-ball is contained in a deterministic finite graph ball,
and every witnessing path is finite.  Monotone convergence proves
the assertion.  The case $M=0$ is immediate.
\end{proof}

\begin{lemma}[A local map in a deleted domain]
\label{mom:map-potential}
Fix one of the maps above, with critical success probability $u>0$,
cap $M\ge2$, and $H>1$.  Run its fixed rule at parameter $p$.  On a
map-dependent interval $p_0\le p<\pc$, with $p_0\ge\pc/2$, its success
probability $u_p$ and conditional endpoint kernel $K_p$ satisfy
\[
 u_p\ge u/2,\qquad \norm{K_p}\le e^{-(H-1)}.
\]
For a deterministic set $A\subseteq V$, put
$F_A=f_{V\setminus A}/\chi$, with value zero on $A$.  There are
constants $Z,C_{\mathrm{map}}$, uniform on this interval, such that
$2\le Z\le CM/u$ and, for every $v\in O$,
\begin{equation}
 (K_pF_A)(v)\le ZF_A(v)+C_{\mathrm{map}}h_A(v),
 \qquad
 h_A(v)=\sum_{\substack{a\in A\\\dist(v,a)\le M}}
                    \sum_{b\sim a}F_A(b).
 \label{eq:s17}
\end{equation}
In the action of $K_p$, functions are restricted to $O$.  The constant
$C$ depends only on the graph.
\end{lemma}

\begin{proof}
The unconditional endpoint kernel has finite range and continuous
entries as a function of $p$: its rule queries only finitely many
edge bits, and its auxiliary randomization has a fixed law.
Covariance makes its row sums constant on $O$.  The absolute
difference of two such kernels is invariant on $O\times O$, so
\eqref{eq:s02} gives equal row and column sums.  The finite range and
the row--column norm bound therefore give operator-norm continuity.
Dividing by the positive scalar $u_p\to u$ proves the assertions
after shortening the interval.  The kernel $K_p$ is stochastic and
invariant, and hence has column sums one, although it need not be
symmetric.  The cap and connection conditions are properties of the
accepted local patterns, so every success at $p$ still has a simple
open path of length at most $M$ to its endpoint.

Choose such a path on each success.  If it avoids $A$, its endpoint
lies in the intrinsic $M$-ball in $V\setminus A$.  Thus
\cref{mom:maximal-potential} bounds its unconditional contribution to
the endpoint potential by
$\bigl(1+(1-p)M/p\bigr)F_A(v)$; this case is empty if $v\in A$.
Endpoints in $A$ contribute zero.  In the remaining case let
$a\sim b$ be the last exit of the chosen path from $A$.  Its terminal
segment from $b$ to the endpoint $y$ lies in $V\setminus A$ and has
length $l\le M$.  Positive association with the event that this
deterministic segment is open yields
\[
 f_{V\setminus A}(b)
 \ge \E_p\bigl[\abs{C_y^{V\setminus A}};
                       \text{the segment is open}\bigr]
 \ge p^l f_{V\setminus A}(y).
\]
Therefore $F_A(y)\le p^{-M}F_A(b)$.  Since
$\dist(v,a)\le M$, summing over all possible last exits bounds this
contribution by $p^{-M}h_A(v)$.  Divide the two bounds by $u_p$.
Because $p\ge\pc/2$ and $u_p\ge u/2$, the first coefficient can be
enlarged to a constant $Z\in[2,CM/u]$, and the second to a finite
constant depending on the fixed map.  This proves \eqref{eq:s17}.
\end{proof}

\begin{proposition}[Critical fractional moments]
\label{mom:critical-moments}
For every $0<\alpha<1/2$,
\begin{equation}
 \sup_{v\in V}\E_{\pc}s_v^\alpha<\infty.
 \label{eq:s21}
\end{equation}
\end{proposition}

\begin{proof}
Assume \eqref{eq:s09}, and define
\[
 D(p)=\chi^{-2}\left[
       \sum_{b\sim o}\E_p(s_os_b;\ C_o\ne C_b)\right].
\]
By \eqref{eq:s04}, $(1-p)\chi'=\chi^2D(p)$.  Our first goal is to
prove, for every $\gamma>0$, that $D(p)\ge c_\gamma\chi^{-\gamma}$
when $p$ is sufficiently close to $\pc$.  We will compare two bounds
on the exterior potential reached by repeated endpoint-map steps:
\eqref{eq:s17} bounds it above in terms of $D(p)$, while the small
norm of $K_p$ bounds the contribution of endpoints in the root's
own cluster and hence leaves a positive exterior contribution.

Fix a map and its interval as in \cref{mom:map-potential}.  For
$x\in O$, expose $A=C_x$ in an independent $p$-field, and use the
possibly unnormalized expectation
\[
 \widehat\E[G(A)]=\E_p\left[\frac{s_x}{\chi}G(C_x)\right].
\]
For each fixed vertex $a$, exposure of the exact cluster $C_a$
leaves fresh exterior edges and hence gives
\begin{align*}
 \widehat\E\left[\one_{\{a\in A\}}
                         \sum_{b\sim a}F_A(b)\right]
 &=\chi^{-2}\sum_{b\sim a}
      \E_p\left[s_a\one_{\{x\in C_a\}}f_{V\setminus C_a}(b)\right]
 \\
 &\le \chi^{-2}\sum_{b\sim a}
                 \E_p(s_as_b;\ C_a\ne C_b)
 \le c_\sigma D(p),
\end{align*}
where $c_\sigma=(\min_i\sigma_i)^{-1}$.  The last inequality follows
because the preceding nonnegative expression is constant on each
vertex type and its type average is $D(p)$.  Bounded degree now gives
$\widehat\E h_A(v)\le C_{\mathrm{map}}D(p)$ uniformly in $v$.
Iterate \eqref{eq:s17}, use stochasticity of $K_p$, and note that
$F_A(x)=0$.  Since $Z\ge2$, the geometric sum gives
\begin{equation}
 \widehat\E(K_p^jF_A)(x)
       \le C_{\mathrm{map}} Z^jD(p),\qquad j\ge1.
 \label{eq:s18}
\end{equation}

We next give a lower bound for the same quantity.  Fresh exterior
exposure identifies it as
\[
 \widehat\E(K_p^jF_A)(x)
 =\chi^{-2}\sum_{v\in O}K_p^j(x,v)
                  \E_p(s_xs_v;\ C_x\ne C_v).
\]
Without the distinctness condition, positive association and the
uniform near-critical comparison of susceptibilities give a lower
bound $c_0>0$ for this normalized sum.  To control what is removed,
write $I=C_x\cap O$, which is nonempty.  Mass transport restricted to
$O$ gives the exact identity
\begin{equation}
 \sum_{v\in O}K_p^j(x,v)\E_p(s_x^2;\ v\in C_x)
 =\E_p\left[\frac{s_x^2}{\abs I}
                  \ip{\one_I}{K_p^j\one_I}\right]
 \le \norm{K_p}^{j}\E_p s_x^2
 \le C\norm{K_p}^{j}\chi^3.
 \label{mom:diagonal-loss}
\end{equation}
For completeness, the transport sends, from $z\in O$ to each
$y\in C_z\cap O$, the mass
\[
 \frac{s_z^2}{\abs{C_z\cap O}}
       \sum_{v\in C_z\cap O}K_p^j(z,v).
\]
Its outgoing and incoming expectations are the two sides of the
identity.  The norm bound uses
$\norm{\one_I}_2^2=\abs I$; it does not require symmetry of $K_p$.
Finally, the three-arm estimate \eqref{eq:s01} gives
$\E_p s_x^2\le\chi_{\max}^3\le C\chi^3$.

Choose $C_*>1$ large enough that
\[
 j=\left\lceil\frac{\log(C_*\chi)}{H-1}\right\rceil
\]
makes the normalized loss in \eqref{mom:diagonal-loss} at most
$c_0/2$.  Combining this lower bound with \eqref{eq:s18} and
$\log Z\le U+m+C'$ gives
\begin{equation}
 D(p)\ge c_{\mathrm{map}}
               \chi^{-(U+m+C')/(H-1)}.
 \label{eq:s19}
\end{equation}
Here $C'$ is graph-dependent; the prefactor may depend on the fixed
map.  Since $U+m\ge\log2$, the maps with $(U+m)/H\to0$ have
$H\to\infty$.  Thus, for any $0<\gamma<1$, we can fix a map for
which the exponent in \eqref{eq:s19} is at most $\gamma$.
The susceptibility identity \eqref{eq:s04} then gives, almost
everywhere on its interval,
$\chi'\ge c\chi^{2-\gamma}$.  In particular
\[
 \frac{d}{dp}\chi^{-(1-\gamma)}
       \le -(1-\gamma)c.
\]
Integrating from $p$ to $r<\pc$, discarding the nonnegative value at
$r$, and then letting $r\uparrow\pc$ yields
\begin{equation}
 \chi(p)\le C_\gamma(\pc-p)^{-1/(1-\gamma)}.
 \label{eq:s20}
\end{equation}
Local absolute continuity, established with \eqref{eq:s04},
justifies this differential calculation.

We transfer this bound to a critical tail by an explicit comparison
of exploration transcripts.  For a large integer $N$, put
$p=\pc-N^{-1/2}$.  From any vertex $v$, explore until the cluster
is exhausted or $N$ vertices have been discovered.  Use a fixed
nonretesting exploration order.  There are at most $dN$ fresh edge
queries, and the transcript decides the event $s_v\ge N$.
If $L_N$ is the likelihood ratio of its law at $\pc$ relative to
its law at $p$, each active query contributes either $\pc/p$ or
$(1-\pc)/(1-p)$.  For fixed $b>1$, its conditional $b$th-moment
factor under $p$ is
\[
 p(\pc/p)^b+(1-p)((1-\pc)/(1-p))^b
       \le 1+C_b(\pc-p)^2.
\]
Indeed this expression is twice continuously differentiable in a
fixed neighborhood of the interior point $\pc$, equals one there,
and has first derivative zero there.  Pad the stopped exploration
with inactive factors equal to one.  Successive conditioning for
the deterministic $dN$ steps gives
\[
 \E_p L_N^b\le(1+C_b/N)^{dN}\le e^{dC_b}.
\]
H\"older's inequality and Markov's inequality therefore imply,
uniformly in $v$,
\begin{equation}
 \Prob_{\pc}(s_v\ge N)
 =\E_p[L_N\one_{\{s_v\ge N\}}]
 \le C_b'\left(\frac{\chi_v(p)}{N}\right)^{1-1/b}.
 \label{mom:transcript-comparison}
\end{equation}
Choose $\gamma$ small and then $b$ large enough that
\[
 \theta=(1-1/b)\left(1-\frac1{2(1-\gamma)}\right)>\alpha,
\]
where $\alpha$ is the exponent in \eqref{eq:s09}.  For all
sufficiently large $N$, \eqref{eq:s20} applies to the chosen $p$;
the uniform type comparison and
\eqref{mom:transcript-comparison} give
$\Prob_{\pc}(s_v\ge N)\le C N^{-\theta}$.
Summing
$((N+1)^\alpha-N^\alpha)\Prob_{\pc}(s_v>N)$ proves a uniform finite
$\alpha$th moment, contrary to \eqref{eq:s09}.  This establishes
\eqref{eq:s21} for every $0<\alpha<1/2$.
\end{proof}

The next argument uses these moments without the temporary
divergence assumption. The conversion from a critical volume tail to
a susceptibility estimate is the principle established in
\cite[Theorem~1.1]{Hutchcroft2020Exponents} for transitive graphs.
We prove the version needed here with constants uniform over the
finitely many vertex types. We include the covariance form of the
O'Donnell--Saks--Schramm--Servedio inequality
\cite[Section~3.3, Equation~(6)]{ODonnellSaksSchrammServedio2005},
in the form needed for all vertex types and all finite domains.

\begin{lemma}[A decision-tree covariance bound]
\label{mom:decision-tree}
Let $X$ have a finite product law, let $0\le g\le1$, and let an
algorithm determine $g(X)$ by querying coordinates without
repetition.  Write $\delta_i$ for the probability that it queries
coordinate $i$, and let $X^{(i)}$ replace just that coordinate by an
independent copy.  For every real function $f$ on the product space,
\begin{equation}
 \abs{\Cov(f(X),g(X))}
 \le\sum_i\delta_i\E\abs{f(X)-f(X^{(i)})}.
 \label{mom:covariance}
\end{equation}
Independent randomization of the algorithm is allowed.
\end{lemma}

\begin{proof}
Take an independent product copy $Y$.  Run the algorithm on $X$,
and let $Z^t$ replace the first $t$ queried coordinates of $X$ by
their values in $Y$.  Include the algorithm's independent seed in
each transcript.  Conditional on any pre-query transcript, the
unqueried coordinates of $X$ retain their independent original
marginals, since that transcript constrains only queried bits.
The replaced coordinates of $Y$ are independent of this transcript
and of those unqueried bits.  Thus the pre-query hybrid has the
original product law conditional on the transcript.  If the next
index is $i$, that index is transcript-measurable and $Y_i$ is still
fresh.  Consequently
\[
 \E\bigl[\abs{f(Z^{t-1})-f(Z^t)}
                     \mid\text{pre-query transcript}\bigr]
       =\E\abs{f(X)-f(X^{(i)})}.
\]
At the terminal transcript, the same reasoning gives product law
for $Z^T$ conditional on the entire transcript.  Since the
transcript determines $g(X)$, it follows that
\[
 \Cov(f(X),g(X))=\E[(f(X)-f(Z^T))g(X)].
\]
Telescope and use $0\le g\le1$ before conditioning on partial
transcripts.  The preceding conditional-increment identity then
proves \eqref{mom:covariance}, because each coordinate is queried
at most once.  Intermediate hybrids need not be independent of
$g(X)$.
\end{proof}

\begin{proposition}[An almost linear susceptibility bound]
\label{mom:susceptibility}
For every $\eta>0$ there is $C_\eta<\infty$ such that
\begin{equation}
 \chi_{\max}(p)\le C_\eta(\pc-p)^{-1-\eta},
 \qquad 0\le p<\pc.
 \label{eq:s22}
\end{equation}
\end{proposition}

\begin{proof}
Fix $0<a<1/2$.  In a finite induced domain, fix a vertex $v$ and
an integer $n\ge2$.  Let $f$ indicate $\abs{C_v}\ge n$.
Independently mark every vertex with probability $1/n$, and let $g$
indicate that $C_v$ contains a mark.  An algorithm deciding $g$
first reveals all marks and then explores the clusters of marked
vertices.  It queries an edge only if at least one endpoint belongs
to a cluster containing a mark.  By \cref{mom:critical-moments},
restriction coupling, and monotonicity in $p\le\pc$, each vertex
has this property with probability at most
\[
 \E_p[1-(1-1/n)^{s_v}]
 \le\E_p\min(1,s_v/n)
 \le n^{-a}\E_p s_v^a\le C_a n^{-a}.
\]
Thus every edge has revealment at most $2C_a n^{-a}$, uniformly in
the domain, the root, and $p\le\pc$.  Mark coordinates have zero
influence on $f$.  For an increasing Boolean edge function, its
absolute resampling influence is exactly
$2p(1-p)$ times the probability that the edge is pivotal.
Writing $F(p)=\Prob_p(s_v\ge n)$, finite-product differentiation
and \cref{mom:decision-tree} give
\[
 \Cov(f,g)\le C_a n^{-a}\,2p(1-p)F'(p).
\]
On the other hand, conditioning on the percolation configuration
gives
\[
 \Cov(f,g)
 \ge\bigl(1-(1-1/n)^n-C_a n^{-a}\bigr)F(p).
\]
For all sufficiently large $n$, the factor in parentheses is
bounded below by a positive constant.  Since $2p(1-p)\le1/2$,
we obtain $F'(p)\ge c_a n^a F(p)$ for $0<p\le\pc$.
Multiplying by $\exp(-c_an^ap)$ shows that
$\exp(-c_an^ap)F(p)$ is nondecreasing.  Integrating up to $\pc$
therefore gives, including the case $F\equiv0$,
\begin{equation}
 \Prob_p(s_v\ge n)
 \le C_a n^{-a}\exp\{-c_a(\pc-p)n^a\}.
 \label{mom:stretched-tail}
\end{equation}
We used the critical moment bound to estimate $F(\pc)$.
Continuity gives the assertion also at $p=0$.  Increasing finite
domain exhaustion is legitimate because the event $s_v\ge n$ has
a finite open witness; all constants are independent of the
domain.  Enlarge the constant to cover the finitely many smaller
values of $n$.

Put $\delta=\pc-p$.  The function
$x^{-a}e^{-c_a\delta x^a}$ is decreasing on $(0,\infty)$.
The substitution $t=c_a\delta x^a$ in its integral gives
\[
 \sum_{n\ge1}n^{-a}e^{-c_a\delta n^a}
 \le C_a\delta^{-(1-a)/a}.
\]
Summing \eqref{mom:stretched-tail} therefore proves the same
bound for $\chi_{\max}$.  Finally take $a=1/(2+\eta)$, for which
$(1-a)/a=1+\eta$, to obtain \eqref{eq:s22}.
\end{proof}

\section{A subpower bound for the triangle trace}
\label{trace:section}

For $p<\pc$ write
\[
 T=T_p,\qquad m_p=\norm{T_p},\qquad a_j=\tr(T^j).
\]
The quantity $a_3$ is the type average of the triangle diagram.
We will prove that it grows more slowly than any prescribed
positive power of $(\pc-p)^{-1}$.

Triangle diagrams and positive connection matrices are classical components of the
tree-graph method of Aizenman and Newman \cite{AizenmanNewman1984};
quantitative relations between triangle diagrams and susceptibility
are developed in \cite{Hutchcroft2022Triangle}. The closed-contact
estimate below is related to the exploration-score and critical-tail
refinements of the two-ghost method
\cite{Hutchcroft2020Locality,Hutchcroft2021PowerLaw}. We prove the
specific contact and trace bounds used here.

For a nonnegative invariant function $F$ of an ordered pair of
distinct finite clusters, define the positive summation functional
\[
 \mathcal D(F)=\left[\E_p\frac1{s_o}
                         \sum_{B\ne C_o}F(C_o,B)\right].
\]
The sum ranges over the other open clusters in the configuration.
For a nonnegative invariant kernel $J$, put
\[
 k_J(A,B)=\sum_{x\in A}\sum_{y\in B}J(x,y),
 \qquad k_e=k_{\Adj}.
\]
In particular $k_e$ counts graph edges between the two clusters.
Mass transport gives the useful rerooting formula
\begin{equation}
 \mathcal D(k_JF)=
 \left[\E_p\sum_v J(o,v)F(C_o,C_v)
                         \one_{\{C_o\ne C_v\}}\right].
 \label{trace:rerooting}
\end{equation}
Indeed send from $x$ to each $y\in C_x$ the mass
$s_x^{-1}\sum_vJ(y,v)F(C_x,C_v)\one_{\{C_x\ne C_v\}}$.
The outgoing and incoming sides of \eqref{eq:s02} are precisely
the two sides of \eqref{trace:rerooting}.  All terms are
nonnegative, so no integrability assumption is needed for this
identity.  Another application, sending
$F(A,B)/(\abs A\abs B)$ from every vertex of $A$ to every vertex
of $B$, shows that $\mathcal D(F(A,B))=\mathcal D(F(B,A))$.

\begin{lemma}[Contacts of two large clusters]
\label{trace:contact-tail}
Fix $0<\alpha<1/2$ and $\epsilon_1>0$ such that
$r=1/2+\epsilon_1-2\alpha<0$.  Uniformly for
$\pc/2\le p<\pc$ and $L\ge1$,
\begin{equation}
 \mathcal D\left(k_e^2
            \one_{\{\min(\abs A,\abs B)\ge L\}}\right)
       \le C L^{1/2+\epsilon_1-2\alpha}.
 \label{eq:s24}
\end{equation}
\end{lemma}

\begin{proof}
Fix an oriented graph edge $(x,v)$ and $n\ge1$.  Let
\[
 H_n=\{C_x\ne C_v,\ n\le\min(s_x,s_v)<2n\}.
\]
On distinct clusters the two size events have disjoint finite open
witnesses.  The BK inequality, coupling with critical percolation,
and \eqref{eq:s21} give
\begin{equation}
 \Prob_p(H_n)\le C n^{-2\alpha}.
 \label{trace:two-large-clusters}
\end{equation}
We also need to control their number of contacts, and for this we
compare two exploration scores.

First explore $C_v$ in the full graph, querying every edge incident
to each discovered vertex exactly once, including edges whose
endpoints have both already been discovered.  Stop after
$N=\lceil2dn\rceil$ queries if the exploration has not finished.
Let $S_{\mathrm{full}}$ be the sum of $\omega_f-p$ over the tested
edges.  Every active query is a fresh Bernoulli($p$) bit.  Padding
after stopping by zero increments, the elementary centered
Bernoulli exponential bound gives
\[
 \E_p e^{\lambda S_{\mathrm{full}}}\le e^{N\lambda^2/8},
 \qquad
 \Prob_p(\abs{S_{\mathrm{full}}}>t)
                       \le2e^{-2t^2/N}.
\]
This is an unconditional estimate, with no cluster-size event
imposed.

On the same percolation configuration, perform a second exploration
by first exposing $A=C_x$ together with its internal bits and
closed boundary.  If $v\notin A$, explore the cluster of $v$ in
$V\setminus A$, again querying every incident edge of that domain
and using the same deterministic query cap $N$.  Let the resulting
score be $S_{\mathrm{ext}}$, and set it to zero if $v\in A$.
Conditional on the exposure, all the queried exterior bits are
fresh and independent.  The identical tail bound therefore holds
conditionally, uniformly in $A$, and hence also unconditionally
after averaging over $A$.

On $H_n\cap\{s_v\le s_x\}$, both explorations finish the same
cluster $B=C_v$, because $\abs B<2n$ and its number of incident
edges is at most $d\abs B<N$.  The completed queried sets are,
respectively, all full-graph edges incident to $B$ and all such
edges with both ends outside $A$.  Their difference consists
exactly of the $A$--$B$ contacts, which are closed.  Query order is
irrelevant to these completed sums, so
\[
 S_{\mathrm{full}}-S_{\mathrm{ext}}=-p k_e(A,B).
\]
It follows that the event under consideration with
$k_e(A,B)>2t/p$ is contained in the union of the two score-tail
events.  Take $t=n^{1/2+\epsilon_1}$, and use $p\ge\pc/2$.
Interchanging $x$ and $v$ handles the opposite size order.  Thus
\[
 \Prob_p\left(H_n\cap
             \{k_e(C_x,C_v)>C n^{1/2+\epsilon_1}\}\right)
             \le C e^{-c n^{2\epsilon_1}}.
\]
Each score was concentrated in its own fresh law before
intersecting with $H_n$; no independence between the scores is
asserted.  Since $k_e\le d\min(s_x,s_v)<2dn$ on $H_n$,
\eqref{trace:two-large-clusters} now yields
\[
 \E_p[k_e(C_x,C_v)\one_{H_n}]
 \le C n^{1/2+\epsilon_1-2\alpha}
                      +C n e^{-c n^{2\epsilon_1}}
 \le C n^r.
\]
Apply \eqref{trace:rerooting} with one contact factor and partition
$\min(\abs A,\abs B)\ge L$ into the bins
$[2^jL,2^{j+1}L)$, $j\ge0$.  There are at most $d$ neighbors of
each root, and $\sum_{j\ge0}(2^jL)^r\le C_rL^r$ because $r<0$.
This proves \eqref{eq:s24}.
\end{proof}

\begin{proposition}[Subpower triangle growth]
\label{trace:subpower}
For every $\eta>0$ there is $C_\eta<\infty$ such that
\begin{equation}
 a_3(p)\le C_\eta(\pc-p)^{-\eta},
 \qquad \pc/2\le p<\pc.
 \label{eq:s23}
\end{equation}
\end{proposition}

\begin{proof}
All operators below are bounded at each subcritical parameter.
The operator-norm local absolute continuity supplied by
\eqref{eq:s04}, and cyclicity of the finite trace, give
$a_3'=3\tr(T^2T')$ almost everywhere.  Set $J=T^2$.
For fixed terminals $x,y$, the first identity in \eqref{eq:s04}
can also be written
\[
 (1-p)T'_p(x,y)
   =\E_p[k_e(C_x,C_y)\one_{\{C_x\ne C_y\}}].
\]
Each edge between the two terminal clusters has exactly one
orientation from the first cluster to the second.  Applying
\eqref{trace:rerooting} to $J$ therefore gives
\begin{equation}
 (1-p)a_3'=3\mathcal D(k_e k_J)
 \quad\text{for almost every }p<\pc.
 \label{eq:s25}
\end{equation}

We will use the following square bound:
\begin{equation}
 \mathcal D(k_J^2)\le\tr(JTJT)=a_6\le m_p^3a_3.
 \label{trace:kernel-square}
\end{equation}
To check the first inequality, reroot one $J$ factor using
\eqref{trace:rerooting} and expand the other one.  This gives
\[
 \left[\sum_{v,x',y'}J(o,v)J(x',y')
   \Prob_p(o\leftrightarrow x',\ v\leftrightarrow y',\
                                      C_o\ne C_v)\right].
\]
The two connections on this event have disjoint open witnesses,
so BK bounds their probability by $T(o,x')T(v,y')$.
Symmetry of $J$ and $T$ identifies the resulting sum with
$\tr(JTJT)$.  Because $J=T^2$, this is $a_6$.  Finally, positivity
of $T$ gives $T^6\le m_p^3T^3$ in quadratic-form order, which
proves the remaining inequality after taking the trace.

Choose $\alpha$ and $\epsilon_1$ as in
\cref{trace:contact-tail}.  Split \eqref{eq:s25} according to
whether the smaller terminal cluster has size less than $L$.
For the part where $A$ is small, reroot one contact and expose
$A=C_o$ with its internal bits and boundary.  If $v\notin A$,
its exterior cluster is fresh, so
\[
 \E_p[k_J(A,C_v)\mid A\text{ and its exposure}]
 \le\sum_{y\in A}\sum_z J(y,z)T(v,z)
 =\sum_{y\in A}(T^3)(y,v).
\]
The positive operator $T^3$ satisfies
\[
 \abs{(T^3)(y,v)}
 \le\sqrt{(T^3)(y,y)(T^3)(v,v)}\le c_\sigma a_3,
\]
since every diagonal entry is bounded by $a_3/\min_i\sigma_i$.
Consequently
\[
 \mathcal D(k_e k_J\one_{\{\abs A<L\}})
 \le d c_\sigma a_3\,[\E_p(s_o\one_{\{s_o<L\}})]
 \le C L^{1-\alpha}a_3.
\]
Here $s\one_{\{s<L\}}\le L^{1-\alpha}s^\alpha$ and
\eqref{eq:s21} supply the last bound.  Reversing the ordered
cluster pair gives the same estimate when $B$ is small, since
$k_e$ and $k_J$ are symmetric in their two arguments.

For the remaining part, Cauchy--Schwarz for the positive measure
underlying $\mathcal D$, followed by \eqref{eq:s24} and
\eqref{trace:kernel-square}, gives
\[
 \mathcal D(k_e k_J\one_{\{\min(\abs A,\abs B)\ge L\}})
 \le C L^{1/4+\epsilon_1/2-\alpha}m_p^{3/2}\sqrt{a_3}.
\]
The factor $(1-p)^{-1}$ in \eqref{eq:s25} is uniformly bounded
on $[\pc/2,\pc)$.  Also $a_3\ge1$, directly from the nonnegative
kernel sum and $T(o,o)=1$.  Thus $b=\sqrt{a_3}\ge1$ is locally
absolutely continuous and satisfies
\begin{equation}
 b'\le C L^{1-\alpha}b
        +C L^{1/4+\epsilon_1/2-\alpha}m_p^{3/2}
 \quad\text{almost everywhere}.
 \label{eq:s26}
\end{equation}

We make the exponent choices explicit.  Given $\eta>0$, choose
$0<\zeta<\min(1,\eta)$, and then choose
$1/2-\alpha>0$, $\epsilon_1>0$, and $\eta'>0$ sufficiently small
that $1/2+\epsilon_1-2\alpha<0$ and
\begin{align*}
 A_*&=(2-\zeta)(1-\alpha)<1,\\
 B_*&=(2-\zeta)(1/4+\epsilon_1/2-\alpha)
                       +\tfrac32(1+\eta')<1+\eta/2.
\end{align*}
These choices are possible by continuity: at the limiting values
$\alpha=1/2$, $\epsilon_1=\eta'=0$, the two expressions are
$1-\zeta/2$ and $1+\zeta/4$, respectively.  Put
$\delta=\pc-p$ and $L=\delta^{-(2-\zeta)}\ge1$.
The susceptibility bound \eqref{eq:s22} implies
$m_p\le\chi_{\max}(p)\le C\delta^{-1-\eta'}$, so
\eqref{eq:s26} becomes
\[
 b'(p)\le C\delta^{-A_*}b(p)+C\delta^{-B_*}.
\]
Since $A_*<1$, the integral of the first coefficient from
$\pc/2$ to $\pc$ is finite.  Multiplication by its integrating
factor therefore gives
\[
 b(p)\le C\left(b(\pc/2)+
                 \int_{\pc/2}^p(\pc-r)^{-B_*}\,dr\right)
       \le C_\eta\delta^{-\eta/2}.
\]
For the last inequality one may bound the integrand by a constant
times $(\pc-r)^{-1-\eta/2}$, using $B_*<1+\eta/2$.
Squaring proves \eqref{eq:s23}.
\end{proof}

\section{Deformation by pivotal corridors}\label{def:section}

We now introduce a deformation of connection probabilities.  It penalizes a
connection whenever that connection must traverse an isolated open corridor.
The penalty is calibrated by a random-walk weight.  Consequently its effect
on the operator norm is small, whereas its effect on the susceptibility can
be estimated from the two ends of the corridor.

Throughout this section fix
\[
 \pc/2\le p<\pc,\qquad q=1-p,\qquad k\in\N,\quad k\ge3.
\]
Only the deformation parameter $t\ge0$ varies.  Write
\[
 T_0=T_p,\qquad m_0=\norm{T_0},\qquad
 \chi_0=\chi(p),\qquad \chi_0^+=\max_x\E_p|C_x|,
 \qquad c_\sigma=(\min_i\sigma_i)^{-1}.
\]
The walk kernel $P$, its norm $\rho<1$, and the type weights $\sigma_i$
are as in \cref{pre:mass-transport,pre:walk}.  All constants in this section depend only on
the graph, unless otherwise indicated.  In particular, their values are
uniform in the displayed range of $p$.

\begin{definition}[Templates and weighted connections]\label{def:templates}
A raw path is a $k$-step path $w=(w_0,\ldots,w_k)$ of $P$, including possible
holding steps.  Its base, endpoint, and probability are
\[
 u=w_0,\qquad v=w_k,\qquad
 \mu_u(w)=\prod_{i=0}^{k-1}P(w_i,w_{i+1}).
\]
Chronological loop erasure, with holding steps removed, gives a simple path
$\gamma_w$ from $u$ to $v$; write $\ell_w$ for its length and $I_w$ for its
internal vertices.  Retain $w$ as a \emph{template} precisely when
$\ell_w\ge3$.  The \emph{blockers} of a template $w$ are all edges incident to
$I_w$ that are not edges of $\gamma_w$; write $b_w$ for their number.  Give
this template the weight
\[
 c(w)=\mu_u(w)p^{-\ell_w}q^{-b_w}.
\]
Different raw paths remain different templates even if they have the same
loop erasure.

For an induced domain $D$, a template is eligible if $\gamma_w\subset D$.
Its raw path need not lie in $D$.  Its weight is always the full-graph weight
$c(w)$ just defined.  A blocker absent from the induced graph on $D$ counts
as closed.

If $x,y\in D$ are connected, let $N^D_{xy}$ be the sum of $c(w)$ over
eligible templates for which every edge of $\gamma_w$ is open and pivotal
for $x\leftrightarrow y$ in $D$, and every blocker is closed or absent.
Here an open edge is pivotal when closing it destroys this connection.
Set $N^D_{xy}=0$ on disconnection, and define
\[
 \begin{split}
 W_t^D(x,y)&=\one_{\{x\leftrightarrow y\text{ in }D\}}
                    e^{-tN^D_{xy}},\\
 X_t^D(x)&=\sum_{y\in D}W_t^D(x,y),\qquad
 s_t^D(x)=\E_p X_t^D(x).
 \end{split}
\]
Connections and potentials with a terminal outside $D$ are zero.
For $x\in D$, $W_t^D(x,x)=1$.  In the full graph put
\[
 T_t(x,y)=\E_pW_t^V(x,y),\qquad
 \chi_t=[s_t^V(o)],\qquad m_t=\norm{T_t}.
\]
\end{definition}

The fixed weights and the convention for absent blockers will allow us to
compare different domains without changing the deformation.

\begin{lemma}[Basic bounds and monotonicity]\label{def:basic}
There is a fixed $C>0$ such that, with $c_k=e^{Ck}$,
\begin{equation}\label{def:edge-weight}
 \sum_{w:\,e\in E(\gamma_w)}c(w)\le c_k/k
 \quad\text{for every edge }e.
\end{equation}
On a connection in a finite cluster,
$N^D_{xy}\le c_k|C_x^D|$.
For fixed $t$, $W_t^D(x,y)$ is increasing under edge openings and under
domain enlargement, using restriction of the same configuration.
The kernel $T_t$ is invariant, symmetric, and entrywise between zero and
$T_0$, and
\begin{equation}\label{def:row-norm}
 m_t\le\sup_x s_t^V(x)\le c_\sigma\chi_t.
\end{equation}
The functions $T_t(x,y)$ and $s_t^D(x)$ are continuously differentiable for
$t\ge0$, with a right derivative at zero, by differentiating
$W_t^D(x,y)$ to $-N^D_{xy}W_t^D(x,y)$ under expectation and summation.
\end{lemma}

\begin{proof}
Since $\ell_w\le k$ and $b_w\le dk$, there is a fixed $B\ge1$ with
$p^{-\ell_w}q^{-b_w}\le B^k$.  For a specified edge $e$, symmetry and
stochasticity of $P$ show that the total probability, summed over all
bases, of traversing $e$ at a specified step is $1/d$.  A union bound over
the $k$ steps therefore gives
\[
 \sum_{w:\,e\in E(\gamma_w)}c(w)
 \le B^k\sum_{u}\sum_{w:\,w\text{ traverses }e}\mu_u(w)
 \le B^k k/d.
\]
Increasing $C$ proves \eqref{def:edge-weight} for every $k\ge3$.
Every counted template has all its path edges on every simple open
$x$--$y$ path.  Charge it to an edge of one such path, which has fewer than
$|C_x^D|$ edges, to obtain the asserted bound on $N^D_{xy}$.

To prove monotonicity, consider a pair already connected before an edge
opening or a domain enlargement.  Every edge of a template counted
afterward belongs to every old simple connection.  Its retained path was
therefore already eligible, open, and pivotal.  Each blocker closed or
absent afterward was closed or absent before.  Thus every template counted
afterward was counted before, with exactly the same weight.  Its total
penalty cannot increase.  A pair previously disconnected has weight zero,
so monotonicity holds for all pairs.

Symmetry in the terminals follows directly from the pivotal condition;
it does not require any reversal property of loop erasure.  Invariance and
entrywise domination are immediate, and the row and column sum bound for
operator norms gives \eqref{def:row-norm}.
Finally,
\[
 \sum_y N^D_{xy}W_t^D(x,y)\le c_k|C_x^D|^2.
\]
The right side has finite expectation by \eqref{eq:s01} and subcritical
susceptibility finiteness.  Dominated convergence proves the differentiation
and continuity assertions.  These arguments apply to infinite domains:
their clusters are almost surely finite, and the preceding bound is
uniform under finite exhaustion.  In fact, once an exhaustion of $D$
contains the entire cluster $C_x^D$, all weighted connections inside it agree with their
limiting values: boundary blockers are closed both before and after their
other endpoints enter the domain.
\end{proof}

\begin{lemma}[Corridor calibration]\label{def:calibration}
For a template $w$, let $Q_w$ be the full-graph event that the edges of
$\gamma_w$ are open and all its blockers are closed.  Put
$D_w=V\setminus I_w$.  Then $Q_w$ is independent of the percolation in
$D_w$, and
\begin{equation}\label{def:calibration-identity}
 c(w)\Prob_p(Q_w)=\mu_u(w).
\end{equation}
On $Q_w$, the template is counted for $x\leftrightarrow y$ if and only if
one of the following two exterior events occurs:
\[
 \begin{split}
 &x\in C_u^{D_w},\quad y\in C_v^{D_w},\quad
                C_u^{D_w}\ne C_v^{D_w},\\
 &x\in C_v^{D_w},\quad y\in C_u^{D_w},\quad
                C_u^{D_w}\ne C_v^{D_w}.
 \end{split}
\]
In particular, neither terminal belongs to $I_w$.
\end{lemma}

\begin{proof}
All bits specified by $Q_w$ are incident to $I_w$, and there are exactly
$\ell_w$ specified open edges and $b_w$ specified closed edges.  This proves
independence and \eqref{def:calibration-identity}.
On $Q_w$ each internal vertex has exactly the two corridor edges open.
The resulting geometry is shown in \cref{fig:corridor}.
A simple connection using every corridor edge must traverse the whole
corridor consecutively, with both terminals outside its interior.  If the
two exterior endpoint clusters coincide, there is an exterior bypass, so
the corridor edges cannot all be pivotal.  If the exterior clusters are
distinct, the corridor is their only open link, so all of its edges are
pivotal for every terminal pair in either displayed orientation.
\end{proof}

\begin{figure}[tbp]
\centering
\begin{tikzpicture}[x=1cm,y=1cm,font=\small,
  dot/.style={circle,fill,inner sep=1.8pt},
  open/.style={very thick,blue!65!black},
  closed/.style={densely dashed,red!65!black}]
\fill[blue!5,rounded corners=12pt] (-1.9,-.95) rectangle (.3,.95);
\draw[blue!30!black,rounded corners=12pt] (-1.9,-.95) rectangle (.3,.95);
\fill[blue!5,rounded corners=12pt] (7.7,-.95) rectangle (9.9,.95);
\draw[blue!30!black,rounded corners=12pt] (7.7,-.95) rectangle (9.9,.95);
\node at (-1,.55) {$C_u^{D_w}$};
\node at (8.9,.55) {$C_v^{D_w}$};
\node[dot,label=below:$x$] (x) at (-1.35,-.35) {};
\node[dot,label=above:$u$] (u) at (0,0) {};
\node[dot,label=above:$v$] (v) at (8,0) {};
\node[dot,label=below:$y$] (y) at (9.2,-.35) {};
\draw[open] (x)--(-.8,-.12)--(u);
\draw[open] (v)--(8.6,-.12)--(y);
\foreach \q in {1.6,3.2,4.8,6.4}{
  \node[dot] at (\q,0) {};
  \draw[closed] (\q,0)--(\q+.25,.7);
  \draw[closed] (\q,0)--(\q-.25,-.7);
  \draw[red!65!black] (\q+.055,.37)--(\q+.195,.29);
  \draw[red!65!black] (\q-.195,-.29)--(\q-.055,-.37);
}
\draw[open] (u)--(v);
\draw[decorate,decoration={brace,amplitude=4pt}] (1.3,1.1)--(6.7,1.1)
  node[midway,above=5pt] {interior $I_w$};
\node[align=center,text width=10.5cm] at (4,-1.5)
{The corridor edges are open; all other edges incident to $I_w$ are closed.};
\end{tikzpicture}
\caption{A forced corridor and its two exterior legs. On the event that
the entire corridor is pivotal for $x\leftrightarrow y$, its endpoints
belong to distinct clusters in $D_w=V\setminus I_w$. The event $Q_w$
fixes only bits incident to $I_w$ and is independent of the exterior
percolation.
Dashed stubs represent closed blocker edges; they are schematic.}
\label{fig:corridor}
\end{figure}


\begin{proposition}[Small change of operator norm]\label{def:norm}
For every $t\ge0$,
\begin{equation}\label{eq:s27}
 0\le T_0-T_t\le 2t\,T_0P^kT_0\quad\text{entrywise},\qquad
 m_t\ge m_0-2t\rho^km_0^2.
\end{equation}
\end{proposition}

\begin{proof}
Expand $-\partial_tT_t(x,y)$ as a nonnegative sum over templates and
discard the factor $e^{-tN^V_{xy}}\le1$.  By
\cref{def:calibration}, the term for $w$ is at most
\[
 \mu_u(w)\bigl(T_0(x,u)T_0(v,y)+T_0(x,v)T_0(u,y)\bigr).
\]
Indeed the two exterior connections in each orientation belong to distinct
clusters, so BK bounds their joint probability by the product of the
corresponding full-graph connection probabilities.  Add the discarded raw
paths to this upper bound and sum over bases and endpoints.  Each orientation
gives $T_0P^kT_0$, since $P^k$ is symmetric.  Integration from zero to $t$
proves the entrywise assertion.  Domination by a nonnegative kernel gives
\[
 \norm{T_0-T_t}\le 2t\norm{T_0P^kT_0}
                  \le2t\rho^km_0^2.
\]
The triangle inequality proves the norm assertion.  In particular, no
positive-semidefiniteness assertion for $T_t$ is used.
\end{proof}

The norm estimate is one half of the comparison. We now state the
corresponding lower bound on the rate at which the averaged row sum
decreases. Let $\E_{\rm path}$ denote averaging over the type-weighted
root $u=o$ and a length-$k$ $P$-path from $u$. Denote its endpoint by
$v$, whether or not the path is retained as a template, and put
$a_3=\tr(T_0^3)$.

\begin{proposition}[Decrease of the averaged susceptibility]\label{def:decrease}
For all $t\ge0$,
\begin{equation}\label{eq:s31}
 -\chi_t'\ge e^{-tc_k}
 \left(
 \E_{\rm path}\left[
       \one_{\{\ell_w\ge3\}}s_t^{D_w}(u)s_t^{D_w}(v)\right]
       -C'(\chi_0^+)^2\rho^ka_3
 \right),
\end{equation}
where $C'>0$ depends only on $G$.
The product term is defined to be zero for discarded raw paths.
\end{proposition}

The product measures the weighted connection mass available at the
two exterior endpoints. We will prove this bound by estimating the
loss under deletion, separating the two exterior clusters, and
comparing their penalties with that of the full connection. To make the
product dominate the error, we still need paths on which both endpoint
potentials remain large. Constructing such paths is the task of the next
section.

\begin{lemma}[Deletion estimate]\label{def:deletion}
For $A\subset D$, let $\Nb[A]$ denote $A$ together with all its neighbors
in the full graph.  For every vertex $v$,
\begin{equation}\label{eq:s28}
 s_t^{D\setminus A}(v)
 \ge s_t^D(v)-\chi_0^+\sum_{z\in\Nb[A]}T_0(v,z).
\end{equation}
\end{lemma}

\begin{proof}
We first compare weighted connections in a fixed configuration.  Suppose
$W_t^D(v,y)>W_t^{D\setminus A}(v,y)$.  There was an old connection.  If
deletion destroys it, an old simple $v$--$y$ path visits $A$; this includes
the case of a deleted terminal.  Otherwise some positive-weight template
is newly counted after deletion.  Its retained path already lay in $D$ and
was open there.  The raw path and its weight have not changed, so there are
only two possible reasons it was not counted before.

First, a path edge may previously have been nonpivotal.  An old simple
connection avoiding that edge must visit $A$, since otherwise the edge
would remain nonpivotal after deletion.  Second, a blocker may previously
have been open.  Since it is now absent, it joins an internal retained-path
vertex $z\notin A$ to a vertex of $A$.  Every surviving simple connection
uses the new template path and therefore visits this
$z\in\Nb[A]$.  This second case includes a removed pendant branch: its
attachment vertex, rather than the removed vertex, lies on the connection.

The entire positive loss, which is at most one, is thus bounded by the
indicator that some old simple $v$--$y$ path visits $\Nb[A]$.  Splitting
such a path at a visited vertex $z$ gives edge-disjoint connection
witnesses from $v$ to $z$ and from $z$ to $y$.  A union bound and BK yield
\[
 \E_p\bigl(W_t^D(v,y)-W_t^{D\setminus A}(v,y)\bigr)
 \le\sum_{z\in\Nb[A]}T_0(v,z)T_0(z,y).
\]
There is no sum over penalty weights in this estimate: a single visiting
event bounds the whole loss.  Summing over $y$ proves \eqref{eq:s28}.
For infinite $A$ all sums are interpreted monotonically, with an infinite
right-hand error giving a trivial bound.
\end{proof}

\begin{lemma}[Two distinct exterior clusters]\label{def:correlation}
There is a constant $C'>0$ such that, for every induced domain $D$ and
$u,v\in D$,
\begin{equation}\label{eq:s29}
 \E_p\bigl(X_t^D(u)X_t^D(v);\ C_u^D\ne C_v^D\bigr)
 \ge s_t^D(u)s_t^D(v)-C'(\chi_0^+)^2T_0^3(u,v).
\end{equation}
The constant is independent of $D,p,k,t$.
\end{lemma}

\begin{proof}
Expose the exact finite cluster $A=C_u^D$, including every internal edge
bit and its closed boundary.  This determines $X_t^D(u)$: every counted
template has retained path in $A$, its weight is deterministic, and its
blockers leaving $A$ are closed or absent.  The edges with both endpoints
in $D\setminus A$ remain independent Bernoulli variables of parameter $p$.

For terminals in a cluster outside $A$, weighted connections before and
after deleting $A$ agree exactly.  All open paths stay in that exterior
cluster; every eligible counted retained path stays there too; the only
blockers removed were closed contacts with $A$.  Raw excursions into $A$
affect neither eligibility nor weight.  Consequently, including $v\in A$
by the zero convention,
\begin{equation}\label{def:cluster-exposure}
 \E_p\bigl(X_t^D(u)X_t^D(v);\ C_u^D\ne C_v^D\bigr)
 =\E_p\bigl[X_t^D(u)s_t^{D\setminus A}(v)\bigr].
\end{equation}
Apply \cref{def:deletion} inside this expectation and use
$X_t^D(u)\le |A|$.  For adjacent $a,z$, positive association with the event
that their edge is open gives $pT_0(v,z)\le T_0(v,a)$.  Hence
\[
 \sum_{z\in\Nb[A]}T_0(v,z)
 \le(1+d/p)\sum_{a\in A}T_0(v,a).
\]
The three-arm estimate \eqref{eq:s01}, summed over a third terminal, gives
\[
 \begin{split}
 \E_p\bigl(|A|\one_{\{a\in A\}}\bigr)
 &=\sum_y\Prob_p(u\leftrightarrow a,\ u\leftrightarrow y\text{ in }D)\\
 &\le\sum_b T_0(u,b)T_0(b,a)\sum_y T_0(b,y)
 \le\chi_0^+T_0^2(u,a).
 \end{split}
\]
The error in \eqref{def:cluster-exposure} is therefore at most
\[
 (1+d/p)(\chi_0^+)^2\sum_a T_0^2(u,a)T_0(a,v),
\]
which proves \eqref{eq:s29} with $C'=1+2d/\pc$.
\end{proof}

\begin{lemma}[Decomposition into two legs]\label{def:legs}
On $Q_w$, suppose
$x\in C_u^{D_w}$, $y\in C_v^{D_w}$, and
$C_u^{D_w}\ne C_v^{D_w}$.  Then
\begin{equation}\label{eq:s30}
 N^V_{xy}\le c_k+N^{D_w}_{xu}+N^{D_w}_{vy}.
\end{equation}
\end{lemma}

\begin{proof}
The total weight of all templates whose retained paths use any edge of
$\gamma_w$ is at most $\ell_wc_k/k\le c_k$, by
\eqref{def:edge-weight}.  Consider any remaining template counted for
$x\leftrightarrow y$.  Its retained path cannot meet $I_w$: every open edge
incident to $I_w$ is an edge of $\gamma_w$.  Its connected retained path
thus lies in one of the two distinct exterior endpoint clusters.

If it lies in the $u$ cluster, every one of its edges is pivotal for the
exterior connection $x\leftrightarrow u$.  Otherwise an exterior
$x$--$u$ path avoiding that edge, followed by the corridor and an exterior
$v$--$y$ path, would give a global $x$--$y$ connection avoiding it.  Its
blockers remain closed or absent on restriction to $D_w$.  Thus it is
counted, with the same weight, in $N^{D_w}_{xu}$.  The $v$ cluster gives
$N^{D_w}_{vy}$ in the same way.  This proves the bound.
\end{proof}

\begin{proof}[Proof of \cref{def:decrease}]
Expand the negative derivative of $\chi_t$ as the sum over the starting
terminal, the other terminal, and all counted templates.  Use the
mass-transport identity \eqref{eq:s02} to move the root from the starting
terminal to the template base.  All terms are nonnegative, so this
rearrangement is justified by Tonelli's theorem.  For each template retain
only the orientation $x\in C_u^{D_w}$, $y\in C_v^{D_w}$.
On $Q_w$ and this exterior event, \cref{def:legs} implies
\[
 W_t^V(x,y)\ge
 e^{-tc_k}W_t^{D_w}(x,u)W_t^{D_w}(v,y).
\]
Independence of $Q_w$ from the exterior and
\eqref{def:calibration-identity} replace its factor $c(w)$ by $\mu_u(w)$.
After summing over $x,y$, we obtain
\[
 -\chi_t'\ge e^{-tc_k}\E_{\rm path}\left[
 \one_{\{\ell_w\ge3\}}
 \E_p\bigl(X_t^{D_w}(u)X_t^{D_w}(v);
           C_u^{D_w}\ne C_v^{D_w}\bigr)\right].
\]
Apply \cref{def:correlation}.  For its nonnegative error kernel, adding
discarded paths only increases the average, and
\[
 \E_{\rm path}T_0^3(u,v)
 =[\sum_v P^k(o,v)T_0^3(v,o)]
 =\tr(P^kT_0^3)\le\rho^k\tr(T_0^3).
\]
The last inequality uses the positive operator $T_0^3$ and the positive
trace inequality of \cref{pre:trace} for the self-adjoint operator $P^k$.  It requires no
commutation of $P$ with $T_0$.  This proves \eqref{eq:s31}.
\end{proof}

The deletion error can be smaller at a vertex near a corridor endpoint
than at the endpoint itself. The following comparison transfers the
potential from that vertex back along a fixed path.

\begin{lemma}[Comparison along a path]\label{def:endpoint}
If a simple path in $D$ joins $b$ to $z$ and has length $l$, then
\begin{equation}\label{eq:s32}
 s_t^D(b)\ge e^{-tlc_k}p^l s_t^D(z).
\end{equation}
\end{lemma}

\begin{proof}
Let $F$ be the event that this fixed path is open.  On $F$ and on a
connection from $z$ to $y$, there is also a connection from $b$ to $y$.
Templates counted for the latter connection whose retained paths use an
edge of the fixed path have total weight at most $lc_k$, by
\eqref{def:edge-weight}.  For any other counted template, every retained
edge is pivotal for $z\leftrightarrow y$: a $z$--$y$ bypass, concatenated
with the fixed $b$--$z$ path, would otherwise give a $b$--$y$ bypass.
Its blocker condition is unchanged.  Thus, on $F$,
\[
 W_t^D(b,y)\ge e^{-tlc_k}W_t^D(z,y).
\]
Both $\one_F$ and $W_t^D(z,y)$ are increasing functions of the percolation
bits, by \cref{def:basic}.  Positive association and $\Prob_p(F)=p^l$
give
\[
 \E_pW_t^D(b,y)
 \ge e^{-tlc_k}\E_p[\one_FW_t^D(z,y)]
 \ge e^{-tlc_k}p^l\E_pW_t^D(z,y).
\]
Summation over $y$ proves \eqref{eq:s32}.
\end{proof}

\section{Paths with an escape at each endpoint}
\label{path:section}

We retain the notation of the corridor deformation. In particular,
$\pc/2\le p<\pc$, $T_0=T_p$, and $I_w$ is the interior of the
chronological loop erasure of a raw $k$-step $P$-path $w$ from $u$ to $v$.
The domain $D_w=V\setminus I_w$ depends only on this path. The next lemma
finds paths for which deleting $I_w$ leaves useful potentials accessible
from both endpoints. All probabilities in the lemma concern auxiliary
walks, independently of the percolation configuration.

Rooted reversibility and mass transport provide the general framework
for the construction; see \cite[Theorem~4.1]{AldousLyons2007}.
The labeled, type-conditioned experiment below supplies the particular
reversal and domination identities needed here. It makes no assertion
that chronological loop erasure commutes with path reversal.

\begin{lemma}[Two endpoint escapes]
\label{path:selection}
For every $\tau>0$ there are constants $J\in\N$, $C_1>0$, $k_0\in\N$,
and $b_*>0$, depending only on $G$ and $\tau$, with the following property.
For every $\pc/2\le p<\pc$, every $k\ge k_0$, and every $R\ge1$, a set
of paths of $\E_{\mathrm{path}}$-probability at least $b_*$ has
loop-erased length at least three and admits, from each of $u,v$, a
simple path in $D_w$ of length at most
\[
 j=J+\left\lceil C_1\log((k+1)R)\right\rceil
\]
to a vertex $z$ satisfying
\begin{equation}
 \sum_{b\in\Nb[I_w]}T_0(z,b)\le\frac{\tau}{R}.
 \label{eq:s33}
\end{equation}
The constants $J,C_1,k_0,b_*$ are fixed before $p,k,R$ are chosen.
\end{lemma}

We first record a smoothing bound which does not involve $\norm{T_0}$.
It is a quasi-transitive, off-diagonal form of the random-walk averaging
estimate known as Schramm's lemma; see \cite[Section~1.1]{Kozma2011}
and \cite[Proposition~6.4]{Hutchcroft2019}. The finite trace gives a
short proof adapted to our symmetric lazy walk.

\begin{lemma}[Smoothing a connection kernel]
\label{path:smoothing}
For all vertices $a,b$ and integers $l\ge0$,
\begin{equation}
 (P^lT_0)(a,b)\le\sqrt{c_\sigma}\,\rho^l.
 \label{eq:s34}
\end{equation}
\end{lemma}

\begin{proof}
Positive-semidefinite Cauchy--Schwarz for $T_0$ gives
\[
 \abs{(P^lT_0)(a,b)}^2
 \le (P^lT_0P^l)(a,a)T_0(b,b).
\]
Here $T_0(b,b)=1$. The invariant positive operator $P^lT_0P^l$ has
every diagonal entry bounded by $c_\sigma$ times its weighted trace.
Consequently, trace cyclicity and the positive-trace bound give
\[
 (P^lT_0P^l)(a,a)
 \le c_\sigma\tr(P^{2l}T_0)
 \le c_\sigma\norm{P^{2l}}\tr(T_0)
 =c_\sigma\rho^{2l}.
\]
This argument uses no commutation between $P$ and $T_0$.
\end{proof}

\begin{proof}[Proof of \cref{path:selection}]
Let $L_{\mathrm{type}}$ and $c_*>0$ be such that every transition
probability of the finite type chain at every time $L\ge
L_{\mathrm{type}}$ is at least $c_*$. The same $c_*$ works for all
types and will therefore remain fixed as $J$ varies.

\paragraph{The labeled experiment.}
For an integer $J$ to be chosen, take a geodesic segment
$\gamma=(\gamma_{-J},\ldots,\gamma_J)$. Such segments exist because an
infinite connected locally finite graph has infinite diameter. Write
$\mathcal S=\Gamma\gamma$ for its orbit as an ordered full segment,
$O_i=\Gamma\gamma_i$, and $\vartheta_i$ for the root weight of $O_i$.
Different indices may describe the same vertex orbit. For $x\in O_i$,
let
\[
 N_i=\#\{g\in\mathcal S:g_i=x\}.
\]
This is a positive finite number independent of the choice of $x$ in
$O_i$. Indeed an anchored segment is contained in the finite ball
$B(x,2J)$, and automorphisms give bijections between anchor sets.
The full orbit $\mathcal S$ itself need not be finite.

Fix $u\in O_0$ and put $L=k-2J\ge L_{\mathrm{type}}$. Sample $g$
uniformly among segments with $g_0=u$. Follow its positive half from
$u$ to $g_J$, then run a middle $P$-walk $\eta$ of length $L$ conditioned
to end in $O_J$. Its conditioning probability
\[
 q_L=P^L(a,O_J),\qquad a\in O_J,
\]
is constant on $O_J$ and is at least $c_*$. Given $\eta$, sample $h$
uniformly among segments with $h_J=\eta_L$, and follow
$(h_J,h_{J-1},\ldots,h_0)$ to $v=h_0$. The resulting raw path $w$ has
length $k$.

Independently, after these choices, take ordinary $P$-walks $\alpha$
and $\beta$ of a common length $m$ from $g_{-J}$ and $h_{-J}$.
Prepend the corresponding negative segment to each walk. These are the
two candidate escapes, of length $J+m$; see
\cref{fig:escape-templates}.
\begin{figure}[tbp]
\centering
\begin{tikzpicture}[x=1cm,y=1cm,font=\small,
  dot/.style={circle,fill,inner sep=1.6pt},
  raw/.style={very thick,blue!65!black,-{Stealth[length=2mm]}},
  escape/.style={thick,orange!70!black,-{Stealth[length=2mm]}}]
\node[dot,label={above left:$u=g_0$}] (u) at (0,0) {};
\node[dot,label=above:$g_J$] (a) at (1.5,1.35) {};
\node[dot,label=above:$h_J$] (b) at (6.5,1.35) {};
\node[dot,label={above right:$v=h_0$}] (v) at (8,0) {};
\node[dot,label=left:$g_{-J}$] (c) at (-1.1,-1.1) {};
\node[dot,label=right:$h_{-J}$] (d) at (9.1,-1.1) {};
\node[dot,label=below:$z$] (z) at (0,-2.65) {};
\node[dot,label=below:$z'$] (zz) at (8,-2.65) {};
\draw[raw] (u)--node[above left] {$J$} (a);
\draw[raw] (a) to[out=-10,in=170] (2.6,1.1)
  to[out=-10,in=180] (4,1.3) to[out=0,in=190] (b);
\node[align=center] at (4,2.65) {middle walk of length $L$\\
conditioned to end in the type of $g_J$};
\draw[raw] (b)--node[above right] {$J$} (v);
\draw[escape] (u)--node[left] {$J$} (c);
\draw[escape] (v)--node[right] {$J$} (d);
\draw[escape] (c) to[out=-80,in=180] (-.3,-1.9) to[out=0,in=180] (.5,-2.1)
  to[out=0,in=0] (z);
\draw[escape] (d) to[out=-100,in=0] (8.3,-1.9) to[out=180,in=0] (7.5,-2.1)
  to[out=180,in=180] (zz);
\node[align=center,text width=5cm] at (4,-1.2)
{Walk tails\\of length $m=j-J$};
\end{tikzpicture}
\caption{The anchored path experiment. The blue raw path has a
prescribed prefix, a reversible middle walk, and a prescribed suffix.
The two orange escapes use the opposite halves of the same geodesic
segments and walk tails that, conditional on $g,\eta,h$, are independent
ordinary $P$-walks before avoidance is imposed. Reversing the middle
walk and exchanging the two segments and escape tails gives the reversed
raw-path law after transport on the center orbit. Avoidance is tested
against the raw path; the escapes may intersect each other. Loop-erasing
either escape keeps it outside the corridor interior.}
\label{fig:escape-templates}
\end{figure}

Write $Q_u$ for the law of the
full labeled object $\xi=(g,\eta,h,\alpha,\beta)$, and $\mu(\zeta)$
for the product of the transition probabilities along any walk $\zeta$.
Its exact probability is
\begin{equation}
 Q_u(\xi)=
 \frac{\mu(\eta)\mu(\alpha)\mu(\beta)}{N_0N_Jq_L}.
 \label{path:labeled-weight}
\end{equation}

\paragraph{Reversal and comparison with the ordinary walk.}
The labeled law serves two purposes. Reversal lets us estimate avoidance
at one endpoint and transfer the estimate to the other. Domination by
the ordinary walk law then converts successful labeled experiments into
a positive mass of the raw paths used in the deformation.

There is a rooted reversal of this experiment:
\[
 \xi^\dagger=(h,\overleftarrow\eta,g,\beta,\alpha),
\]
rooted at $v$. Symmetry of $P$ shows that
$Q_u(\xi)=Q_v(\xi^\dagger)$ by \eqref{path:labeled-weight}; neither
$N_0=N_J$ nor equality of the center and endpoint type weights is
needed. The escape tails are swapped, rather than reversed. For an
invariant nonnegative function $F$ of the labeled object, apply
\eqref{eq:s02} to
\[
 M_F(x,y)=\sum_{\xi:x\to y}Q_x(\xi)F(\xi),
 \qquad x,y\in O_0,
\]
and set this transport to zero off $O_0\times O_0$. Invariance makes
the outgoing sum constant on $O_0$, and the common root weight
$\vartheta_0$ cancels. The reversal bijection identifies the incoming
sum with the outgoing sum for $F\circ\dagger$. Thus
\begin{equation}
 \E_{Q_u}F=\E_{Q_u}(F\circ\dagger).
 \label{path:rooted-reversal}
\end{equation}

We also need domination after forgetting all labels. For a fixed raw
path $w$, let $n_g$ and $n_h$ count the full segments compatible with
its visible first and last blocks, respectively. Its middle is fixed,
and summing all compatible full segments and all escape tails gives
\[
 Q_u(w)=\frac{\mu(\eta)}{q_L}
        \frac{n_g}{N_0}\frac{n_h}{N_J}
 \le \frac{\mu(\eta)}{c_*}.
\]
The two fractions are at most one. In particular, different negative
halves which project to the same visible half cause no additional
multiplicity. The ordinary walk probability is
$\mu_u(w)=(2d)^{-2J}\mu(\eta)$, so
\begin{equation}
 Q_u(w)\le D_J\mu_u(w),\qquad
 D_J=\frac{(2d)^{2J}}{c_*}.
 \label{path:domination}
\end{equation}

\paragraph{Avoiding the raw path.}
We bound all collisions uniformly in the escape-tail length $m$.
First consider the escape from $u$, conditional on $g$. Its deterministic
negative half and the positive half of $g$ meet only at $u$, which is
allowed. Before conditioning the middle endpoint type, the probability
that middle time $r$ visits any one of the $J+1$ negative-segment
vertices is at most $\rho^r$, and is zero for $r<J$. Similarly, at
escape-tail time $l$ the probability of visiting any one of the $J+1$
positive-segment vertices is at most $\rho^l$, and is zero for $l<J$.
For a collision of the random escape tail and the middle at times
$l,r$, their independence and symmetry of $P$ give
\[
 \sum_xP^l(g_{-J},x)P^r(g_J,x)
 =P^{l+r}(g_{-J},g_J)\le\rho^{l+r}.
\]
This probability is zero when $l+r<2J$. Conditioning the middle
endpoint costs at most $c_*^{-1}$ in these estimates. Hence all these
collisions have total probability at most
\begin{equation}
 E_J=\frac1{c_*}\left(
 2(J+1)\sum_{r\ge J}\rho^r
 +\sum_{n\ge2J}(n+1)\rho^n\right),
 \qquad E_J\longrightarrow0.
 \label{path:early-collisions}
\end{equation}
Extending the time sums to infinity is what makes this bound uniform
in $L$ and $m$.

For collisions with the last block, define the anchored-position kernel
\[
 H_{a,b}(x,y)=
 \one_{\{x\in O_a\}}
 \frac{\#\{g\in\mathcal S:g_a=x,\ g_b=y\}}{N_a}.
\]
Its row sums are one on $O_a$ and zero elsewhere. Its column sum is
constant on $O_b$, is zero elsewhere, and, by \eqref{eq:s02}, is exactly
\begin{equation}
 \sum_xH_{a,b}(x,y)=\frac{\vartheta_a}{\vartheta_b}
 \quad(y\in O_b),\qquad
 \norm{H_{a,b}}\le
 \sqrt{\frac{\vartheta_a}{\vartheta_b}}\le\sqrt{c_\sigma}.
 \label{path:position-kernel}
\end{equation}
Thus the operator bound is independent of $J$ and of all anchor counts.
Conditional on $g$, the middle endpoint distribution $\nu$ satisfies
$\norm{\nu}_2\le c_*^{-1}\rho^L$. Each suffix-position distribution
$\nu H_{J,i}$, $0\le i\le J$, therefore has norm at most
$\sqrt{c_\sigma}\,c_*^{-1}\rho^L$.

Conditional on $g$, the first escape is independent of the middle and
$h$. Its deterministic positions are point masses, and the distribution
at tail time $l$ has norm at most $\rho^l$. Cauchy--Schwarz and a union
bound now show that its collisions with the suffix have probability at
most
\begin{equation}
 F_J\rho^L,\qquad
 F_J=C\left((J+1)^2+\frac{J+1}{1-\rho}\right),
 \label{path:suffix-collisions}
\end{equation}
where $C$ depends only on $G$. The estimate includes all tail times
because $\sum_{l\ge0}\rho^l=(1-\rho)^{-1}$. The same suffix-position
bound at $i=0$ gives
\[
 Q_u(\dist(u,v)<3)
 \le \abs{B(u,2)}^{1/2}\sqrt{c_\sigma}\,c_*^{-1}\rho^L
 \le C_0\rho^L,
\]
with $C_0$ independent of $J$. By \eqref{path:rooted-reversal}, the
second escape has exactly the same avoidance-failure bound as the
first: the relevant event refers only to the raw path and its reversed
raw path.

Choose $J$ so large that $2E_J\le1/8$, and then choose
$L_0\ge L_{\mathrm{type}}$ so large that
$(2F_J+C_0)\rho^{L_0}\le1/8$. For every $L\ge L_0$ and every $m$,
with probability at least $3/4$ both candidate escapes avoid every raw
path vertex other than their respective starting vertices, and
$\dist(u,v)\ge3$. The escapes need not avoid each other. This choice
order is valid because no anchor-count factor occurs in
\eqref{path:early-collisions}.

Avoidance will keep the escape paths inside $D_w$. It remains to choose
their endpoints so that the deletion error in \cref{def:deletion} is
small. For this we use the smoothing bound before imposing avoidance.

\paragraph{Testing the potentials.}
Set $\lambda=-\log\rho>0$, choose $C_1$ with $\lambda C_1>1$, and put
$m=\lceil C_1\log((k+1)R)\rceil$. Conditional on $g,\eta,h$, and
\emph{before imposing any avoidance event}, the two tails are independent
ordinary $P$-walks. Their terminal vertices $Z_1,Z_2$ therefore satisfy,
by \cref{path:smoothing} and $\abs{\Nb[I_w]}\le(d+1)(k+1)$,
\[
 \E\left[\sum_{b\in\Nb[I_w]}T_0(Z_i,b)\,\middle|\,g,\eta,h\right]
 \le C(k+1)\rho^m,\qquad i=1,2.
\]
Here the loop-erased interior is defined whether or not its length is
at least three. Markov's inequality bounds the sum of the two potential
failure probabilities by
\begin{equation}
 2C\tau^{-1}R(k+1)\rho^m
 \le 2C\tau^{-1}((k+1)R)^{1-\lambda C_1}.
 \label{path:potential-failures}
\end{equation}
Increase a fixed $k_0\ge2J+L_0$ until the right-hand side is at most
$1/4$ for all $k\ge k_0$ and $R\ge1$. All constants remain independent
of $p$ and $R$.

Subtracting the failure probability in \eqref{path:potential-failures}
from the preceding $3/4$ success bound leaves probability at least
$1/2$ for both avoidance and both potential tests simultaneously.
Since $\dist(u,v)\ge3$, the retained corridor has length at least three.
Its interior is contained in the raw path and excludes $u,v$. Each
escape therefore lies in $D_w$, and erasing its loops gives a simple
path to the same terminal vertex of length at most $J+m$. No reversal
property of chronological loop erasure is used.

Project this successful event onto raw paths admitting these witnesses.
By \eqref{path:domination}, ordinary walk law from a root in $O_0$
assigns these paths probability at least $1/(2D_J)$. Averaging over
root types proves the lemma with
$b_*=\vartheta_0/(2D_J)>0$.
\end{proof}

\section{Critical boundedness and the nonuniqueness interval}
\label{con:section}

We combine the preceding estimates to prove the critical operator bound.
The argument fixes all path constants before choosing the subcritical
parameter and the deformation time. We then transfer the bound to an
interval above criticality and prove the simultaneous cluster assertion.

\begin{proof}[Proof of \cref{thm:critical-operator}]
Suppose that the critical connection kernel $T_{\pc}$ is unbounded.
Put $p_0=\pc/2$. For $p_0\le p<\pc$, use the corridor notation
\[
 \delta=\pc-p,\qquad S=\log(1/\delta),\qquad
 m_0=\norm{T_p},\qquad R=\frac{\chi_0^+}{m_0}\ge1,
 \qquad a_3=\tr(T_p^3).
\]
Sprinkling \eqref{eq:s05} with $a=\delta/(1-p)$ would bound
$T_{\pc}$ if $ad m_0<1$. The assumed unboundedness, together with
\eqref{eq:s22} and \eqref{eq:s23}, therefore gives
\begin{equation}
 m_0\ge\frac{1-p}{d\delta},\qquad
 \log R=o(S),\qquad \log a_3=o(S)
 \quad(\delta\downarrow0).
 \label{eq:s35}
\end{equation}
Indeed $m_0\ge c\delta^{-1}$, whereas, for every $\eta>0$,
$\chi_0^+\le C_\eta\delta^{-1-\eta}$ and
$1\le a_3\le C_\eta\delta^{-\eta}$. Hence
$1\le R\le C'_\eta\delta^{-\eta}$ for every $\eta>0$.

\begin{claim}[Comparison of deformed row sums]
\label{con:row-comparison}
There is $c_f>0$, depending only on $G$, such that
\begin{equation}
 s_t^V(z)\ge c_f\chi_t
 \qquad\text{for every $z\in V$ whenever $tc_k\le1$}.
 \label{eq:s36}
\end{equation}
\end{claim}

\begin{proof}
Connectedness and the finite number of vertex types give an integer
$L_f$ such that every vertex is within distance $L_f$ of every vertex
orbit. For fixed $t$, choose an orbit on which $s_t^V$ is maximal and
a simple path of length at most $L_f$ from $z$ to that orbit.
Equation \eqref{eq:s32}, with $p\ge p_0$, gives
\[
 s_t^V(z)\ge e^{-L_f}p_0^{L_f}\max_x s_t^V(x)
 \ge e^{-L_f}p_0^{L_f}\chi_t.
\]
Thus $c_f=e^{-L_f}p_0^{L_f}$ is admissible.
\end{proof}

Fix $0<\tau<c_f/(4c_\sigma)$ and the constants of
\cref{path:selection} for this value of $\tau$. They are now fixed
throughout the limit $p\uparrow\pc$. Since
$j\le J+1+C_1\log((k+1)R)$, there are fixed constants $c>0$ and
$C_2>0$ such that
\begin{equation}
 p^{2j}\ge c(kR)^{-C_2}\qquad(k\ge k_0,\ R\ge1).
 \label{eq:s37}
\end{equation}
For example, one may take
$C_2=2C_1\log(1/p_0)$ and
$c=p_0^{2(J+1)}2^{-C_2}$.

Put $\lambda=-\log\rho>0$ and choose a fixed $K_*$ with
$\kappa:=\lambda K_*>C_2+4$. Set
\[
 A_\delta=\log(2R)+\log(2+a_3)+\log S,\qquad
 k=\lceil K_*A_\delta\rceil,\qquad
 B_*=(kR)^{C_2+2},\qquad t_*=\frac{B_*}{m_0}.
\]
By \eqref{eq:s35}, $k\to\infty$ and $k=o(S)$; in particular
$k\ge k_0$ and $k\le S$ eventually. The estimates needed below are
\begin{equation}
 B_*\rho^k=o(1),\qquad
 R^2a_3\rho^k=o((kR)^{-C_2}),\qquad
 t_*jc_k=o(1).
 \label{eq:s38}
\end{equation}
To check their simultaneous validity, the definition of $k$ gives
\[
 \rho^k\le (2R)^{-\kappa}(2+a_3)^{-\kappa}S^{-\kappa}.
\]
Since $R,a_3\ge1$ and $k\le S$, it follows that
\[
 B_*\rho^k\le C S^{C_2+2-\kappa}\longrightarrow0,
 \qquad
 (kR)^{C_2}R^2a_3\rho^k
 \le C S^{C_2-\kappa}\longrightarrow0.
\]
Finally $c_k=e^{Ck}$ and $m_0\ge c e^S$, so
\[
 \log(t_*jc_k)
 \le -S+(C_2+2)(\log k+\log R)+\log j+Ck+O(1)
 =-S+o(S).
\]
This proves \eqref{eq:s38} with every path constant already fixed.

Consider the deformation with these values of $p,k$, on
$0\le t\le t_*$. Equation \eqref{eq:s27} and the first estimate in
\eqref{eq:s38} give
\begin{equation}
 m_t\ge m_0(1-2B_*\rho^k)\ge\frac{m_0}{2},\qquad
 \chi_t\ge\frac{m_0}{2c_\sigma}.
 \label{con:retained-norm}
\end{equation}
The second inequality uses the row-sum bound $m_t\le c_\sigma\chi_t$;
it requires no positivity of $T_t$ as an operator.

For a selected raw path in \cref{path:selection}, let $z$ be either
of its two escape endpoints. By \eqref{eq:s28}, \eqref{eq:s33},
\eqref{eq:s36}, and \eqref{con:retained-norm},
\[
 \begin{split}
 s_t^{D_w}(z)
 &\ge c_f\chi_t-\chi_0^+\frac{\tau}{R}
  =c_f\chi_t-\tau m_0\\
 &\ge(c_f-2c_\sigma\tau)\chi_t
  \ge\frac{c_f}{2}\chi_t.
 \end{split}
\]
The hypotheses of \eqref{eq:s36} hold throughout this interval by the
last estimate in \eqref{eq:s38}. Apply \eqref{eq:s32} to the two simple
escape paths in $D_w$, each of length at most $j$. Their product is
bounded below by
\[
 s_t^{D_w}(u)s_t^{D_w}(v)
 \ge e^{-2tjc_k}p^{2j}\left(\frac{c_f}{2}\chi_t\right)^2.
\]
Such paths have $\E_{\mathrm{path}}$-probability at least $b_*$, so
\eqref{eq:s37} and \eqref{eq:s38} imply
\begin{equation}
 \E_{\mathrm{path}}\!\left[
  \one_{\{\ell_w\ge3\}}s_t^{D_w}(u)s_t^{D_w}(v)\right]
 \ge c_1(kR)^{-C_2}\chi_t^2,
 \label{con:positive-term}
\end{equation}
uniformly on $[0,t_*]$, with $c_1>0$ independent of sufficiently small
$\delta$. On the other hand, the error in \eqref{eq:s31} is at most
\[
 C'(\chi_0^+)^2\rho^ka_3
 =C'R^2m_0^2\rho^ka_3
 \le4C'c_\sigma^2R^2\rho^ka_3\chi_t^2
 =o((kR)^{-C_2}\chi_t^2).
\]
The prefactor $e^{-tc_k}$ in that equation is uniformly bounded below.
We conclude that, for a constant $c_2>0$,
\begin{equation}
 -\chi_t'\ge c_2(kR)^{-C_2}\chi_t^2
 \qquad(0\le t\le t_*).
 \label{con:differential-decrease}
\end{equation}
The differentiability and domination established for the deformation
justify integration. Since $\chi_t>0$, it yields
\[
 \frac1{\chi_{t_*}}
 \ge c_2(kR)^{-C_2}t_*
 =\frac{c_2(kR)^2}{m_0}.
\]
But \eqref{con:retained-norm} gives
$1/\chi_{t_*}\le2c_\sigma/m_0$. These inequalities are incompatible
with $kR\to\infty$. Thus $T_{\pc}$ is bounded, proving
\cref{thm:critical-operator}.
\end{proof}

The next elementary endpoint lemma records the equivalence between
critical boundedness and the uniform subcritical formulation in
\cref{thm:main}.

\begin{lemma}[Monotone passage to an endpoint]\label{con:critical-limit}
Let $G$ be an infinite locally finite graph, let $p_*\in(0,1]$, and let
$T_p(x,y)=\Prob_p(x\leftrightarrow y)$ be its Bernoulli bond connection
kernels. Then $T_{p_*}$ is bounded on counting-measure $\ell^2(V)$ if
and only if $\sup_{0\le p<p_*}\norm{T_p}_{2\to2}<\infty$. In that case
the supremum equals $\norm{T_{p_*}}_{2\to2}$.
\end{lemma}

\begin{proof}
Each finite-domain connection probability is continuous in $p$, and
connectivity in $G$ is the increasing union of finite-domain
connectivity events. Thus $T_p(x,y)\uparrow T_{p_*}(x,y)$ as
$p\uparrow p_*$. If the subcritical norms are bounded by $M$, then
for finitely supported $f,g$ the finite sum
$\ip f{T_pg}$ tends to $\ip f{T_{p_*}g}$ and has absolute value at
most $M\norm f_2\norm g_2$. The limiting sesquilinear form extends
by density to a bounded operator with the required matrix entries
and norm at most $M$. Conversely, entrywise nonnegativity gives
\[
 |\ip f{T_pg}|\le\ip{|f|}{T_p|g|}
 \le\ip{|f|}{T_{p_*}|g|}
 \le\norm{T_{p_*}}\norm f_2\norm g_2.
\]
The two inequalities prove both equivalence and equality of norms.
\end{proof}

The simultaneous phase theorem of H\"aggstr\"om, Peres, and Schonmann
\cite[Theorem~1.3]{HaggstromPeresSchonmann1999} applies to all
quasi-transitive graphs and gives one event valid throughout the open
nonuniqueness phase. The following direct argument proves the compact
interval consequence from operator boundedness alone. It avoids taking
an uncountable intersection of events of probability one.

\begin{proposition}[One coupling for an interval]
\label{con:coupled-interval}
Let $G$ be infinite, connected, locally finite, and quasi-transitive.
Suppose $\pc<p_1<p_2<1$ and $T_{p_2}$ is bounded on $\ell^2(V)$.
In the coupling obtained by assigning independent uniform labels to
edges and opening an edge at parameter $p$ when its label is at most
$p$, almost surely there are infinitely many infinite clusters
simultaneously at every $p\in[p_1,p_2]$.
\end{proposition}

\begin{proof}
Let $N_{12}$ be the number of $p_2$-clusters containing an infinite
$p_1$-cluster. First, every automorphism-invariant event of the label
field has probability zero or one. Indeed, approximate such an event
$E$, within error $\varepsilon$ in probability, by a finite-coordinate
event $A$. Its coordinates lie in a finite ball centered at a vertex
in an unbounded orbit; such an orbit exists because there are finitely
many orbits and infinitely many vertices. An automorphism can move
this ball disjointly from itself. The events $A$ and its translate are
then independent, whereas $E$ agrees with its translate. It follows
that $\abs{\Prob(E)-\Prob(E)^2}\le4\varepsilon$, and letting
$\varepsilon\downarrow0$ proves the assertion.
In particular the invariant, measurable random variable
$N_{12}\in\{0,1,2,\ldots,\infty\}$ is constant almost surely.
Since $p_1>\pc$, an infinite $p_1$-cluster exists with positive
probability, so this constant is positive.

It cannot equal one. Write
$\theta(x)=\Prob_{p_1}(|C_x|=\infty)$. Some $\theta(x)$ is positive;
positive association with a finite path makes $\theta$ positive at
every vertex. Invariance and the finite number of types give
$\theta_*:=\inf_x\theta(x)>0$. If $N_{12}=1$, any two vertices in
infinite $p_1$-clusters belong to the same $p_2$-cluster. Positive
association therefore implies
\[
 T_{p_2}(x,y)\ge\theta(x)\theta(y)\ge\theta_*^2
 \qquad(x,y\in V).
\]
For finite sets $F$ of arbitrarily large size this gives
$\ip{\one_F}{T_{p_2}\one_F}/\norm{\one_F}_2^2
\ge\theta_*^2|F|$, contradicting boundedness.

Suppose instead that $N_{12}=n$ almost surely for an integer $n\ge2$.
By countability there are fixed vertices $x,y$ which, with positive
probability, belong to infinite $p_1$-clusters lying in distinct
$p_2$-clusters. Fix a finite path from $x$ to $y$. Replace its edge
labels by independent uniform labels in $[0,p_1]$. The resulting law
is absolutely continuous with respect to the original product law:
only finitely many coordinates were resampled, each with a bounded
density. At both parameters this modification only opens edges, and
on the specified positive-probability event it merges two counted
$p_2$-clusters.

Every infinite $p_1$-cluster after the modification contains an
infinite $p_1$-cluster from before it. Indeed finitely many added edges
can combine only finitely many old clusters into any new cluster; a
finite union of finite clusters is finite. Thus each new counted
$p_2$-cluster contains an old counted $p_2$-cluster, and old
$p_2$-clusters only merge. On the event in question the new count is
at most $n-1$. This contradicts absolute continuity and the almost
sure equality $N_{12}=n$. Hence $N_{12}=\infty$ almost surely.

On this single event, choose an infinite $p_1$-cluster in each of
infinitely many distinct counted $p_2$-clusters. For every
$p\in[p_1,p_2]$, their containing $p$-clusters are infinite and remain
distinct because their $p_2$-clusters are distinct. This proves the
simultaneous assertion.
\end{proof}

\begin{proof}[Proof of \cref{thm:main}]
First suppose $G$ is simple. If $\Aut(G)$ is unimodular,
\cref{thm:critical-operator} gives $\norm{T_{\pc}}<\infty$.
Since $\pc<1$, sprinkling \eqref{eq:s05} at $p=\pc$ with any
sufficiently small $a>0$ gives a bounded connection kernel at
$\pc+(1-\pc)a>\pc$. Thus $\pc<\ptwo$.

If $\Aut(G)$ is nonunimodular, \cref{pre:unimodular-reduction} supplies both
critical boundedness and the strict operator gap. In either case,
\cref{pre:threshold-order} yields $\pc<\ptwo\le\pu$.

If loops or parallel edges are allowed, use the simple subdivision
$H$ from \cref{pre:simple-reduction}. Infinite-cluster counts in $H$
at parameter $t$ agree with those in $G$ at parameter $t^2$, so
$\pc(H)^2=\pc(G)$. Moreover, on the old vertex set,
\[
 T^H_t\big|_{V(G)\times V(G)}=T^G_{t^2}.
\]
Compression of a bounded operator to this coordinate subspace is
bounded. Applying the simple-graph result to $H$ therefore gives
boundedness of $T^G_{\pc(G)}$ and of $T^G_p$ at some $p>\pc(G)$.
The inequality $\ptwo(G)\le\pu(G)$ again follows from
\cref{pre:threshold-order}. This covers the full graph class in the
statement.

Finally choose $\pc<p_1<p_2<1$ with $T_{p_2}$ bounded, which the strict
operator gap permits. \Cref{con:coupled-interval} supplies
the claimed simultaneous nonempty interval of infinitely many
infinite clusters. The equality of the uniform subcritical and critical
norms is \cref{con:critical-limit}, completing the proof.
\end{proof}

\section{Consequences of the operator theorem}
\label{cons:section}

We now prove Corollary~\ref{cor:critical-laws}. The two parts of the
connection-operator conclusion play related but distinct roles.
Boundedness at criticality makes the triangle diagram finite, bringing
the standard mean-field theory into play.
The strict interval of operator boundedness above $\pc$ gives
exponential two-point decay there and supplies the stronger input
used for the extrinsic-radius law.

\begin{proof}[Proof of Corollary~\ref{cor:critical-laws}]
For every graph in the corollary, \cref{thm:main} makes
$T_{\pc}$ a bounded operator on counting-measure $\ell^2(V)$.
Tonelli's theorem identifies its nonnegative matrix product with the
triangle sum. If $\delta_v$ is the unit vector at $v$, then
\[
 \nabla_{\pc}(v)
 =\ip{\delta_v}{T_{\pc}^{3}\delta_v}
 \le\norm{T_{\pc}^{3}}_{2\to2}
 \le\norm{T_{\pc}}_{2\to2}^{3}.
\]
This proves \eqref{eq:critical-triangle}. The same implication is
recorded in \cite[Section~1]{Hutchcroft2020L2}.

\paragraph{Critical laws on simple graphs.}
First suppose that $G$ is simple. The susceptibility,
percolation-probability, and cluster-size estimates are the classical
mean-field consequences of the triangle condition
\cite{AizenmanNewman1984,BarskyAizenman1991}; their quasi-transitive
form is recorded in
\cite[Section~1, Equations~(1.1)--(1.3)]{Hutchcroft2020L2}.
For explicit applications with the precise quasi-transitive hypotheses,
we use the following later results. Under the stronger condition
$\pc<\ptwo$ available here, we may apply
\cite[Theorem~1.1 and Corollary~1.4]{Hutchcroft2022FiniteClusters}.
The corollary with $k=1$ and $p\uparrow\pc$ gives
\eqref{eq:meanfield-chi}, since subcritical clusters are finite.
The theorem at $p=\pc$ gives \eqref{eq:meanfield-size}, since
Lemma~\ref{pre:critical-finiteness} excludes infinite critical clusters.

For percolation probability, Lemma~2.1 of
\cite{Hutchcroft2022FiniteClusters}, whose hypothesis is
the triangle condition, bounds the intrinsic-radius tail uniformly
over vertices. Letting the radius tend to infinity gives
$\theta_v(p)\le C(p-\pc)$ for $p\ge\pc$: local finiteness makes an
infinite cluster equivalent to an unbounded intrinsic radius.
For the lower bound, choose one representative from each vertex
orbit, forming a finite set $F$. Theorem~2 and Propositions~11
and~23 of \cite{AntunovicVeselic2008}, with
$\beta=-\log(1-p)$, give
\[
 \sum_{u\in F}\theta_u(p)\ge c(p-\pc)
\]
for $p>\pc$ sufficiently close to $\pc$. Indeed, their bound is
linear in $\beta-\beta_c$, and this is comparable to $p-\pc$ near
$0<\pc<1$. Fix $v$ and choose a path of length $l_u$ from $v$ to
each $u\in F$. Harris association applied to the event that this
path is open and the event $\{|C_u|=\infty\}$ gives
$\theta_v(p)\ge p^{l_u}\theta_u(p)$. Hence
\[
 \theta_v(p)\ge
 \frac{\min_{u\in F}\pc^{l_u}}{|F|}
 \sum_{u\in F}\theta_u(p)
 \ge c_v(p-\pc),
\]
which proves \eqref{eq:meanfield-theta}.

For the radius estimates, it is important to distinguish the two
metrics. The intrinsic-radius upper bound under the triangle condition
originates in the lattice setting in
\cite[Theorem~1.2(ii)]{KozmaNachmias2009}; its
quasi-transitive form and the matching general lower bound are
described in \cite[Section~3]{Hutchcroft2020L2}. The extrinsic-radius
law uses the stronger condition $\pc<\ptwo$
\cite[Theorem~3.1]{Hutchcroft2020L2}. Both estimates in the exact
quasi-transitive form needed here are included in
\cite[Theorem~1.2]{Hutchcroft2022FiniteClusters}, whose hypothesis is
$\pc<\ptwo$. Set $p=\pc$ in that theorem. Its events require the
cluster radius to be finite; Lemma~\ref{pre:critical-finiteness} shows that
all critical clusters are finite almost surely, so this restriction
may be removed. The resulting estimates are
\eqref{eq:meanfield-intrinsic} and \eqref{eq:meanfield-extrinsic}.

\paragraph{Exponential connectivity.}
The definition of $\pexp$ above is equivalent to that in
\cite[Section~2]{Hutchcroft2020L2}: a positive exponential rate is
equivalent to an estimate with some positive exponent and finite
prefactor, uniformly over vertices. Theorem~2.2 of that paper gives
$\ptwo\le\pexp$. More precisely, its Remark~2.3 gives, for
$0<p<\ptwo$,
\[
 T_p(x,y)\le
 2\norm{T_p}_{2\to2}
 \exp\left(-\frac{\dist_G(x,y)}
                 {e\norm{T_p}_{2\to2}}\right)
 \qquad(x,y\in V).
\]
At $p=0$ the off-diagonal connection probabilities vanish. Together
with $\pc<\ptwo$ from \cref{thm:main}, this proves
\eqref{eq:exponential-threshold} and
\eqref{eq:exponential-connectivity} for simple graphs.

\paragraph{Subdivision and transfer of the laws.}
It remains to pass from simple graphs to the full graph class of the
corollary. Delete the loops, which change none of the observables,
and let $H=\widetilde G$ be the simple subdivided graph in
Lemma~\ref{pre:simple-reduction}. Denote the set of nonloop edge copies of
$G$ by $E_0$, so $V(H)=V(G)\sqcup E_0$. Let $d_0$ bound the number
of incident nonloop edges at a vertex of $G$, counting parallel copies.
In the coupling of that lemma, parameter $t$ on $H$ induces parameter
$p=t^2$ on $G$, old-vertex connection probabilities agree, and
$\pc(H)^2=\pc(G)$.

For an old root $v$, write $C_v^G$ and $C_v^H$ for the two coupled
clusters, and use the same superscripts for the connection and
percolation observables. Their old vertices agree. Every new vertex
in $C_v^H$ is joined by an open edge to an old neighbor in the same
cluster, and every old vertex has at most $d_0$ new neighbors.
Hence, pathwise,
\begin{equation}\label{eq:subdivision-size}
 |C_v^G|\le |C_v^H|\le(1+d_0)|C_v^G|,
 \qquad
 \theta_v^G(t^2)=\theta_v^H(t).
\end{equation}
For two old vertices, both the ambient distance and the distance along
open paths in $H$ are twice their counterparts in $G$. A new vertex
in the root cluster is at distance at most one from an old vertex in
that cluster. Consequently, for either radius,
\begin{equation}\label{eq:subdivision-radius}
 2\operatorname{rad}_{*}(C_v^G)
 \le \operatorname{rad}_{*}(C_v^H)
 \le2\operatorname{rad}_{*}(C_v^G)+1,
 \qquad *\in\{\mathrm{int},\mathrm{ext}\}.
\end{equation}
The statements include infinite values. Applied at
$t=\pc(H)=\sqrt{\pc(G)}$, these comparisons transfer the three
critical tails from $H$ to $G$. They also transfer susceptibility
and percolation probability on the appropriate sides of criticality:
if $t_c=\sqrt{\pc(G)}>0$, then
$p-\pc(G)=(t-t_c)(t+t_c)$, whose second factor stays bounded above
and away from zero near $t_c$. Multiplicative constants may change,
as allowed in the corollary.

\paragraph{Transfer of the operator threshold.}
To transfer the full range in \eqref{eq:exponential-connectivity},
we also compare the two operator thresholds. Define the incidence
operator $B:\ell^2(E_0)\to\ell^2(V(G))$ by
$B(x,e)=\one_{\{x\text{ is an endpoint of }e\}}$.
Its row sums are at most $d_0$ and its column sums are two.
Cauchy--Schwarz at each old vertex gives
$\norm{Bf}_2^2\le2d_0\norm f_2^2$, so
$\norm B_{2\to2}\le\sqrt{2d_0}$. The operator
\[
 Q:\ell^2(V(G))\longrightarrow
       \ell^2(V(G))\oplus\ell^2(E_0),\qquad
 Qf=(f,B^*f),
\]
satisfies $\norm Q_{2\to2}^{2}\le1+2d_0$.
For an old-to-new connection in $H$, a simple open path must enter the new
vertex through one of its two old endpoints. For a connection between
distinct new vertices, choose an old endpoint at each end of a simple open
path. Between these old endpoints the path induces a connection in
$G$. The union bound, together with equality for old-to-old
connections, therefore gives the entrywise inequality
\[
 T_t^H\le Q T_{t^2}^G Q^*.
\]
The product on the right is interpreted entrywise here; each entry
is a sum over at most four pairs of old endpoints.
This inequality also covers a new diagonal entry: the corresponding
diagonal entry of $B^*T_{t^2}^G B$ is at least two.
If $T_{t^2}^G$ is bounded, entrywise domination of nonnegative
kernels makes $T_t^H$ bounded with norm at most
$\norm Q^2\norm{T_{t^2}^G}$. Conversely, if $T_t^H$ is bounded,
compression to the old vertices makes $T_{t^2}^G$ bounded.
Thus boundedness is equivalent and, in that case,
\[
 \norm{T_{t^2}^G}_{2\to2}
 \le\norm{T_t^H}_{2\to2}
 \le(1+2d_0)\norm{T_{t^2}^G}_{2\to2}.
\]
It follows that
\[
 \ptwo(H)^2=\ptwo(G).
\]
For every $0\le p<\ptwo(G)$, apply the simple-graph exponential
bound to $H$ at $t=\sqrt p<\ptwo(H)$. Old-vertex connectivity agrees
and $\dist_H(x,y)=2\dist_G(x,y)$ for $x,y\in V(G)$, so restriction
gives \eqref{eq:exponential-connectivity} with constants depending
only on $G$ and $p$. Its validity for every $p<\ptwo(G)$ implies
$\ptwo(G)\le\pexp(G)$, completing the proof.
\end{proof}

\Needspace{12\baselineskip}
\begingroup
\small
\bibliographystyle{plain}
\bibliography{references}
\endgroup
\end{document}
